\documentclass[11pt]{article}

%==============================================================================
% Packages
%==============================================================================
\usepackage[T1]{fontenc}
\usepackage[utf8]{inputenc}
\usepackage{lmodern}
\usepackage[margin=1in]{geometry}
\usepackage{amsmath,amssymb}
\usepackage{booktabs,tabularx,array,longtable}
\usepackage{enumitem}
\usepackage{xcolor}
\usepackage[colorlinks=true,linkcolor=blue!50!black,urlcolor=blue!50!black,citecolor=blue!50!black]{hyperref}
\usepackage{url}

%==============================================================================
% Macros
%==============================================================================
\newcommand{\R}{\mathbb{R}}
\newcommand{\E}{\mathbb{E}}
\newcommand{\sg}{\operatorname{sg}}
\newcommand{\RMSNorm}{\operatorname{RMSNorm}}
\newcommand{\expm}{\operatorname{expm}}
\newcommand{\bpb}{\textsc{bpb}}
\newcommand{\code}[1]{\texttt{#1}}
\newcommand{\vram}{\mathrm{VRAM}}
\newcommand{\erank}{\operatorname{erank}}
\newcommand{\MI}{I}
\newcommand{\Entropy}{H}
\newcolumntype{L}[1]{>{\raggedright\arraybackslash}p{#1}}
\definecolor{gate}{rgb}{0.07,0.40,0.18}
\definecolor{diag}{rgb}{0.55,0.40,0.00}
\newcommand{\GATE}{\textcolor{gate}{\textbf{GATE}}}
\newcommand{\DIAG}{\textcolor{diag}{\textbf{DIAG}}}

%==============================================================================
% Title
%==============================================================================
\title{
OPG, Final Minimal Design\\
\large Necessary-and-Sufficient Models, Metrics, Assumptions, and Protocol for P1
}
\author{Consolidated revision for \code{train\_gpt.py}}
\date{2026-06-04}

\begin{document}
\maketitle
\tableofcontents

%==============================================================================
\section{Goal and Claim Hierarchy}
\label{sec:goal}

The project should first establish the smallest decisive claim: under a fixed memory budget,
recurrent depth buys task utility on data that actually requires depth.  The claims are ordered so
that later mechanisms cannot substitute for a failed primary result:
\begin{itemize}[leftmargin=1.4em,itemsep=1pt]
  \item \textbf{P1, primary:} \(G_T(K_{\rm lo},K_{\rm hi})>0\) at fixed memory on a depth-hard
  task, with a passing positive control.
  \item \textbf{P2, memory:} reversible reconstruction makes activation memory approximately
  constant in recurrent depth \(K\).  This claim is largely severable from P1 and must be measured
  independently.
  \item \textbf{P3, secondary:} an MoE basis recovers the depth diversity lost by weight tying.
  This claim is staged only after P1 holds.
\end{itemize}

The design serves P1.  P2 verifies the training-memory niche.  P3 is optional capacity recovery, not
the premise of the first experiment.  The core \emph{retains} the mechanisms that directly advance
the two design goals: MLA (low-rank joint KV compression), smooth-sparse MoE (no hard top-\(k\)), and
a mixture-of-softmaxes (MoS) output head are retained for the resource (MLA low-rank KV; smooth-sparse
MoE compute) and expressiveness (MoE effective-depth basis, MoS output-rank) goals.  What it removes
are feature stacks that obscure P1 failure attribution rather than advancing the goals: Dirichlet--UCB
routing, implicit/DEQ-style backward, fixed-point consistency, spectral or Lyapunov pressure, UFID,
distillation, and anytime losses are excluded from the P1 science path.  The executable
pipeline is only the positive-control, \(M_0\), and \(M_{\rm clk}\) triage.  MoE, cache, and
deployment axes are executable only through a separate staged diagnostic harness.  This preserves
the causal split: S+0 promotes or fails P1, while S+1/S+2/S+3 report mechanism and deployment
diagnostics without becoming explanations for a failed P1 gate.

%==============================================================================
\section{Models: The Necessary Set}
\label{sec:models}

\paragraph{End-to-end \(M_0\) flow (forward + backward).}
The implemented pipeline (\code{train\_gpt.py}) is, in order:
\begin{enumerate}[leftmargin=1.6em,itemsep=1pt]
  \item \textbf{Embed.}  \(x_0=\code{\detokenize{tok_emb}}(\mathit{tokens})+\code{\detokenize{pos_emb}}\);
  the embedding is tied and reused as the MoS output basis.
  \item \textbf{Reversible recurrence over a sampled depth \(K\)} (additive coupling, two streams):
  \(a_{k+1}=Q\,a_k+F_\theta(\RMSNorm(b_k+x_0+e_k))\),
  \(b_{k+1}=Q\,b_k+G_\theta(\RMSNorm(a_{k+1}+x_0+e_k))\), where \(F_\theta,G_\theta\) are pre-norm delta
  blocks (MLA attention + SwiGLU MoE), \(e_k\) is the optional step-clock (\(M_{\rm clk}\),
  \code{step\_conditioning}, off by default), and \(Q\) is the optional orthogonal (matrix-exponential
  of a low-rank skew) damping (\code{--damping}, default \code{none}\(\Rightarrow Q{=}I\)): strictly
  invertible (\(Q^{\top}\)), norm-preserving (no fixed-point collapse, extrapolates beyond the trained
  depth), and relaxing the \(A{+}B{=}I\) convex tie.
  \item \textbf{Reversible \(O(1)\)-in-\(K\) backward} (\code{RevRecurrenceFn}): reconstructs activations
  from the terminal state via the exact inverse, with fp64 stream accumulation
  (\code{\detokenize{recurrence_accum_dtype}}) and replayed bf16 autocast keeping
  \(\code{\detokenize{recon_rel}}\approx0\) so the BPTT gradients are correct under bf16.
  \item \textbf{Readout, outside the recurrence.}  \(z_K=\tfrac12(a_K+b_K)\to\code{\detokenize{final_norm}}\)
  (RMSNorm) \(\to\) MoS head (\(\code{\detokenize{n_mix}}\) mixture-of-softmaxes, log-probs) \(\to\) nll.
  \item \textbf{MoE health.}  Default-on DeepSeek loss-free bias balancing + ST-MoE z-loss; per-token
  entropy-max demoted to off; optional architecture-search routing knobs.
  \item \textbf{Finite-horizon training.}  \(K_{\rm hi}\sim\code{\detokenize{k_set}}\), shallow
  \(k_{\rm lo}\); loss \(=\) final-depth nll \(+\) a stop-grad no-degradation hinge (Eq.~\eqref{eq:p1-loss}),
  \emph{not} a fixed-point objective.
\end{enumerate}
The subsections below give the precise forms.

\paragraph{Core \(M_0\).}
The core model is a tied, token-injected, additive-coupling reversible recurrence whose \(F_\theta\)
and \(G_\theta\) blocks are pre-norm MLA attention~\cite{deepseekv2} plus a smooth-sparse SwiGLU
mixture-of-experts (MoE).  The recurrent state is
\(s_k=(a_k,b_k)\), with input embedding state \(x_0\):
\begin{align}
  a_{k+1}
  &=
  a_k+F_\theta\!\left(\RMSNorm(b_k+x_0)\right),
  \\
  b_{k+1}
  &=
  b_k+G_\theta\!\left(\RMSNorm(a_{k+1}+x_0)\right).
\end{align}
The inverse is structural and floor-free:
\begin{align}
  b_k
  &=
  b_{k+1}-G_\theta\!\left(\RMSNorm(a_{k+1}+x_0)\right),
  \\
  a_k
  &=
  a_{k+1}-F_\theta\!\left(\RMSNorm(b_k+x_0)\right).
\end{align}
Invertibility is therefore an algebraic property of the additive coupling, not a consequence of
contraction.  The core does not require a fixed point.  Under mixed-precision (bf16 autocast) training
the coupling streams \((a,b)\) are accumulated in fp64 (the additive \(\pm\) uses fp64 add/sub) while
\(F_\theta,G_\theta\) keep their matmuls under the ambient bf16 autocast; the backward replays that
autocast so reconstruction recomputes the identical bf16 \(F_\theta,G_\theta\).  This makes the inverse
bit-exact (\(\code{\detokenize{recon_rel}}\approx0\)) and the \(O(1)\)-memory reversible BPTT gradients
correct under bf16 (\S\ref{sec:metrics}, the \code{\detokenize{recurrence_accum_dtype}} knob).

The readout is simple: for a recurrent budget \(K\), the final recurrent state
\[
  z_K=\tfrac12\!\left(a_K+b_K\right)
\]
is normalized by a single readout RMSNorm, \(z=\RMSNorm(z_K)\), and \(z\) is fed to a
mixture-of-softmaxes (MoS) output head~\cite{mos}, which mixes a small number of softmax
components over the tied embedding basis to break the softmax-rank bottleneck on the output
distribution.  The readout norm is applied strictly \emph{outside} the reversible recurrence, so
exact (floor-free) reconstruction is unaffected; it is required because \(z_K\) accumulates \(K\)
additive updates from \(x_0\) and grows with depth, and an unnormalized \(z_K\) saturates the
MoS head's \(\tanh\) context map, zeroing its gradient.  There is no halting gate: depth is selected by the sampled/evaluated recurrent budget
\(K\), and the recurrence's \(F_\theta\)/\(G_\theta\) blocks use MLA attention and smooth-sparse
SwiGLU MoE, trained with full reversible BPTT~\cite{revffn,revnet} (\(O(1)\) activation memory in
\(K\), \S\ref{sec:related}) so deeper inference is not penalized by activation storage.  The MoE router is a control-experiment
knob, \texttt{router\_type}\(\in\{\texttt{relu},\texttt{softmax}\}\) (ReLU routing per
ReMoE~\cite{remoe} is continuous, reconstruction-stable, and reversibility-safe; cf.\
MoEUT~\cite{moeut}).

\paragraph{Configurable recurrence-block structure (architecture-search axes).}
The \(F_\theta\)/\(G_\theta\) delta block exposes a family of composable, reversibility-preserving
structural knobs so a controlled architecture search can run without touching the recurrence or its
floor-free inverse.  Every axis defaults to the core \(M_0\) behavior, and each keeps the block a
deterministic pure function of its single input (smooth routing only; no top-\(k\)/argmax/capacity
dispatch), so exact fp64 reconstruction holds for every setting.
(1) \emph{Block order} \(\in\{\texttt{attn\_ffn},\texttt{ffn\_attn},\texttt{parallel}\}\) selects how
attention and the FFN-MoE compose inside one delta sub-block:
\(\texttt{attn\_ffn}\) (default) computes \(a=\mathrm{attn}(\RMSNorm(\mathit{inp}))\) then
\(a+\mathrm{moe}(\RMSNorm(\mathit{inp}+a))\); \(\texttt{ffn\_attn}\) reverses the two;
\(\texttt{parallel}\) sums both branches from the same \(\RMSNorm(\mathit{inp})\).
(2) \emph{Attention MoE} (off by default) replaces the single shared MLA with a MoEUT-style~\cite{moeut}
per-token MoE over \(n_{\rm attn}\) independent low-rank MLA experts (each owning its Q/KV
compression--decompression, K-rope, and output projection), combined with the same smooth-routing
family and trained through the same router auxiliaries.
(3) \emph{Shared experts} \(n_{\rm shared}\ge0\) add DeepSeek-style~\cite{deepseekv2} always-on,
ungated experts summed into every token in addition to the routed experts,
\(\mathrm{moe}(x)=\sum_s\mathrm{shared}_s(x)+\sum_e r_{t,e}\,\mathrm{routed}_e(x)\), applied to the FFN
MoE (and the attention MoE when enabled).
(4) \emph{Expert layout} \(n_{\rm sub}\ge1\) makes each \(F_\theta\)/\(G_\theta\) map a stack of
\(n_{\rm sub}\) unique attn+MoE sub-blocks applied in sequence as one clean pure delta
(\(h\leftarrow\mathit{inp}\); \(h\leftarrow h+\mathrm{sub}(h)\) per sub-block; return \(h-\mathit{inp}\)),
trading experts-per-sublayer against unique-sublayers at matched parameter/compute
(e.g.\ \(16\times1\) vs \(8\times2\)).
(5) \emph{Expert B-init} \(\in\{\texttt{small},\texttt{zero}\}\) sets the routed-expert output-projection
init: \(\texttt{small}\) is the non-zero engagement default; \(\texttt{zero}\) is the classic LoRA-B
init, only sensible with \(n_{\rm shared}>0\) providing a live base.

\paragraph{Collapse prevention: loss-free balancing + z-loss (default).}
The active-expert fraction is no longer a resource metric (the resource goal is
activation memory; \S\ref{sec:metrics}), so dense routing is acceptable and the only objective is to
keep experts \emph{alive and diverse} (expressiveness / effective depth).  At \(\approx\)10\,M scale
the softmax router does \emph{not} suffer global dead-expert collapse (global utilization
\(\code{\detokenize{expert_util}}\) stays high) but instead per-token \emph{over-confidence}:
per-token router entropy \(\code{\detokenize{router_entropy}}\) crashes (\(2.4\to0.06\), each token
near one-hot).  The entropy-MAX regularizer (\S below) MAXIMIZES per-token entropy --- an antagonistic
objective that fights useful specialization --- and did not prevent this collapse, so it is
\emph{demoted to off-by-default}.  The default collapse preventers are now two related-work mechanisms
that are gradient-interference-free and reversibility-safe:
\begin{enumerate}
\item \emph{DeepSeek-V3 auxiliary-loss-free load balancing}~\cite{deepseekv3} (default, the SOTA
  balancer):  a per-expert bias \(b\in\mathbb{R}^E\) is ADDED to the router logits before the softmax,
  \(r=\mathrm{softmax}(\code{router}(x)+b)\), and updated once per optimizer step by the sign rule
  \(b_e\!\leftarrow\!b_e + \eta_b\,\mathrm{sign}(\bar L - L_e)\) (raise under-used, lower over-used),
  then re-centered (\(b\!\leftarrow\!b-\bar b\)), where \(L_e\) is the mean routing weight of expert
  \(e\) accumulated over tokens AND recurrence steps and \(\bar L\) its mean.  Crucially \(b\) is a
  \emph{detached buffer} (no gradient, in no optimizer group), updated ONLY at the step boundary, so it
  is CONSTANT within any single forward+backward --- the reversible inverse recomputes the identical
  routing map, leaving reconstruction (\(\code{\detokenize{recon_rel}}\approx0\)) and grad-equivalence
  exact.  It is the permitted ``slow bias feedback'': global usage pressure with no auxiliary-loss
  gradient interference.  The accumulated load is all-reduced across ranks before the update so every
  rank keeps an identical bias.  Knob \(\texttt{router\_bias\_update\_rate}\) (\(\eta_b\), default
  \(10^{-3}\); 0 disables).
\item \emph{ST-MoE router z-loss}~\cite{stmoe} (default; the direct anti-over-confidence term):  the
  loss adds \(c_z\,Z\) with \(Z=\frac{1}{BT}\sum_{t}\big(\mathrm{logsumexp}_e\,\ell_{t,e}\big)^2\) over
  routed tokens (\(\ell\) the router logits).  It bounds the router logit magnitude, preventing the
  runaway confidence that drove per-token entropy to \(0.06\).  It is an in-graph router-only term
  (recomputed on the saved block input), so reversibility is untouched.  Knob
  \(\texttt{router\_z\_coef}\) (\(c_z\), default \(10^{-3}\); 0 disables).
\end{enumerate}
The forward routing map is otherwise unchanged.  Both default mechanisms reach the router only because
the experts are \emph{not} inert: the per-expert output projection \(W^{\rm out}_e\) is small-non-zero
initialized (the attention output projection alone is zero-initialized for readout stability), so each
expert contributes a small non-zero delta from step zero and both the experts and the router receive a
non-zero task gradient.  A zero-initialized \(W^{\rm out}_e\) would zero every expert output, so the
router weights would multiply a zero and receive \emph{no} task gradient --- the MoE would be inert.

\paragraph{Demoted / ablation auxiliaries (control-experiment axis).}
The following remain composable CLI knobs for explicit ablation but are \emph{off by default} (a
\emph{fixed} \(L_1\) penalty on ReLU routing weights, by contrast, monotonically drives every weight to
zero, \(\code{\detokenize{active_frac}}\to0\)):
\begin{enumerate}
\item \emph{Entropy regularization} (demoted, off by default; maximize): the loss \emph{subtracts}
  \(c_H\,H(p)\) where \(H(p)=-\sum_e p_e\log p_e\) is the mean-token entropy of the router
  distribution (softmax weights, or ReLU weights renormalized to a distribution with zero-mass tokens
  skipped).  Higher entropy \(\Rightarrow\) more uniform expert use.  Demoted because it MAXIMIZES
  per-token entropy --- fighting useful specialization --- and did not prevent the measured collapse.
\item \emph{Switch-Transformer load balance}~\cite{switch} (minimize): the loss adds
  \(c_{\rm sw}\,E\sum_e f^{\rm top}_e\,P_e\), with \(P_e=\frac{1}{BT}\sum_t r_{t,e}\) the mean router
  probability mass (differentiable) and \(f^{\rm top}_e\) the fraction of tokens whose \(\arg\max\)
  (top) expert is \(e\) (detached; the \(\arg\max\) is used \emph{only} in this loss factor, never in
  the forward dispatch, so reversibility holds).  Minimal (\(=1\)) at a uniform load, maximal (\(=E\))
  when one expert carries everything.
\item \emph{ReMoE adaptive-\(\lambda\) controller}~\cite{remoe} (ReLU router, enabled by
  \(\texttt{router\_target\_active\_frac}>0\)): a scalar \(\lambda_{\rm route}\) is updated once per
  optimizer step toward a target sparsity \(S^\star=1-\texttt{router\_target\_active\_frac}\),
\[
  \lambda_{\rm route}\ \leftarrow\ \mathrm{clamp}\!\Big(\lambda_{\rm route}\cdot\alpha^{\,\mathrm{sign}(S_{\rm meas}-S^\star)},\ 10^{-8},\ 10^{3}\Big),\qquad \alpha=1.2,
\]
  raising \(\lambda_{\rm route}\) when routing is too dense and lowering it when too sparse, which
  holds \(\code{\detokenize{active_frac}}\) near the target rather than collapsing.  Its aux penalty is the
  load-balanced \(L_1\) form \(\mathcal{L}_{\rm lb}=\frac{1}{E}\sum_e f_e\,\bar r_e\) with per-expert usage
  fraction \(f_e=\frac{1}{BT}\sum_{t}\mathbf{1}[r_{t,e}>0]\) and mean mass \(\bar r_e=\frac{1}{BT}\sum_t r_{t,e}\),
  so over-used experts are penalized more.  The realized sparsity
  \(S_{\rm meas}\) is measured once per optimizer step at the controller boundary (all-reduced
  across ranks for a rank-consistent \(\lambda_{\rm route}\)), never inside the gradient-accumulation
  micro loop.  Softmax routing is dense (\(S_{\rm meas}\approx0\)), so the controller is a no-op.
\end{enumerate}
Diagnostics logged: per-token router entropy
(\(\code{\detokenize{router_entropy}}\), specialization), global expert-utilization entropy
(\(\code{\detokenize{expert_util}}\), utilization), and \(\code{\detokenize{diag:active_frac}}\).
Two diagnostics directly probe the MoE-basis-depth mechanism (experts composed across recurrence steps
\(\Rightarrow\) exponentially many depth-\(K\) expert paths --- what matters is whether routing
\emph{varies across steps}):  \(\code{\detokenize{route_step_div}}\) is the mean over tokens of the
number of DISTINCT dominant (\(\arg\max\)) experts across the \(K\) recurrence steps divided by \(K\)
(\(1.0\): a different expert every step, maximal exponential-path diversity; \(1/K\): the same expert
every step, degenerate);  \(\code{\detokenize{expert_cos_div}}\) is the mean pairwise \(1-\cos\) between
the experts' output vectors on a probe batch (higher: a more diverse basis; \(\approx0\):
degenerate/duplicated experts).  Both are log-site reads (no hot-loop sync).
An effective-depth diagnostic \(\code{\detokenize{disp_tail}}\) is also logged: the tail-mean over the
last half of the depth steps of the per-step relative recurrence displacement
\(\lVert z_{k+1}-z_k\rVert/\lVert z_k\rVert\) (with \(z_k=\tfrac12(a_k+b_k)\) the reversible midpoint).
A sustained (non-decaying) displacement signals the recurrence keeps doing work at depth (high
effective depth); a rapid decay to \(\approx0\) signals early saturation.
The reversible round-trip error \(\code{\detokenize{recon_rel}}\) is also logged as the BPTT
\emph{gradient-correctness gate}: the relative error of inverting the terminal state
\((a_K,b_K)\) back to the seed \((a_0,b_0)\),
\(\code{\detokenize{recon_rel}}=(\lVert a_0'-a_0\rVert+\lVert b_0'-b_0\rVert)/(\lVert a_0\rVert+\lVert b_0\rVert+\epsilon)\).
The \(O(1)\)-memory reversible backward reconstructs activations from \((a_K,b_K)\) via exactly this
inverse, so gradients equal ordinary BPTT only if the round-trip is exact.  Mixed-precision
correctness rests on two design choices.  \emph{(i) fp64 stream accumulation.}  The recurrence carries
the coupling streams \((a,b)\) in fp64 (the \code{\detokenize{recurrence_accum_dtype}} knob, default
\code{float64}; \code{float32} is a cheaper option), so the additive coupling \(\pm\) is exact while
the \(F_\theta,G_\theta\) matmuls still run at the ambient autocast precision (bf16); their output is
cast to the stream dtype before each \(\pm\).  This keeps the algebraic inverse bit-exact and removes
the catastrophic low-precision \(\pm\) drift that otherwise amplifies over \(K\) steps in
ill-conditioned (random-init) regimes.  \emph{(ii) Autocast replay in the backward.}  The forward runs
\(F_\theta,G_\theta\) under the caller's bf16 autocast; the backward \emph{captures and replays} that
autocast so the reconstruction and the rebuilt-step VJP recompute \(F_\theta,G_\theta\) identically
(same bf16 matmuls) --- otherwise an fp32 backward would invert a \emph{different} function than the
bf16 forward and corrupt the gradient.  With both, \(\code{\detokenize{recon_rel}}\approx0\) (machine
precision) under bf16 autocast, so the reconstructed graph --- and the gradient signal --- match the
true forward; the readout (\code{\detokenize{final_norm}}/\code{\detokenize{mos_head}}/loss) is cast
back to the model dtype so only the recurrence stream is fp64.  A large
\(\code{\detokenize{recon_rel}}\) means the reconstructed graph drifts from the true forward.  It is a
log-site-only probe (one forward sweep plus one inverse sweep), measured under the training bf16
autocast so it reports the actual training reconstruction drift.
The trained model can still fail; the point is that the architecture composes
shared low-rank experts per step into high \emph{effective} depth.  Nestedness across depth is a
property of the additive-coupling recurrence and budget selection, not of a halting gate.

The memory-critical target implementation uses reversible full BPTT through \(K_{\rm hi}\).
Truncation is only a calibrated fallback, because full reversible BPTT is the core advantage for
P2.  The Tier-1 executable harness may use ordinary stored-activation autograd for detectability
and structural-inverse tests; it must not claim P2 activation-memory scaling until the
memory-saving backward path is implemented and measured.  The P1 loss is:
\begin{equation}
  \mathcal L
  =
  \E_{K_{\rm hi}\sim q_{\rm deep}}\!\left[L_T(K_{\rm hi})\right]
  +
  \lambda_h
  \left[
    L_T(K_{\rm hi})-\sg(L_T(K_{\rm lo}))+m
  \right]_+,
  \qquad
  \lambda_h\ \text{small}.
  \label{eq:p1-loss}
\end{equation}
Hard \(\mathrm{NaN}/\mathrm{Inf}\) guards and a light boundedness penalty,
\[
  [\|z_k\|_{\rm RMS}-B_z]_+^2,
\]
are stability guards, not fixed-point losses.

\emph{Implementation note (same-start hinge under random init).}  When
\code{--init-state random} is active, \code{M0GPT.forward} draws a fresh random
recurrence seed \((a_0,b_0)\) per call.  The \code{train\_gpt.py} finite-horizon
path therefore captures and \emph{resets} the RNG before each of the two passes
\(L_T(K_{\rm lo})\) and \(L_T(K_{\rm hi})\), so both start from the \emph{same}
seed noise and the no-degradation hinge compares depth budgets \(K_{\rm lo}\) vs
\(K_{\rm hi}\) from an identical start (a controlled same-start depth comparison;
a no-op for \code{--init-state x0}, which draws no noise).

\emph{Implementation note (Tier-1 harness).}  The executable Tier-1 harness
(\code{experiments/p1\_synthetic.py}) currently optimizes only the
\(L_T(K_{\rm hi})\) cross-entropy term; it does \emph{not} yet apply the
\(\lambda_h\) no-degradation hinge or the boundedness penalty above.  Those terms
are realized in the \code{train\_gpt.py} finite-horizon path.  The S+0 detectability
results are therefore from the plain-\(L_T\) objective.

\paragraph{Clock diagnostic \(M_{\rm clk}\).}
\(M_{\rm clk}\) is \(M_0\) with an explicit step embedding \(e_k\) injected into
\emph{both} half-step updates each iteration:
\[
  a_{k+1}=a_k+F_\theta\!\left(\RMSNorm(b_k+x_0+e_k)\right),\qquad
  b_{k+1}=b_k+G_\theta\!\left(\RMSNorm(a_{k+1}+x_0+e_k)\right),
\]
where \(e_k\) is sinusoidal or learned.  This model is the upper bound on depth-conditioned
computation.  It is diagnostic only and excluded from the final model unless its extrapolation cost
is measured and accepted.

\paragraph{Positive control.}
The harness must include a vanilla looped-transformer positive control: input-injected, recurrent
depth exposed, activations stored rather than reconstructed, and otherwise deliberately ordinary.
If this model does not scale on the Tier-1 depth-hard task, the harness, task, or scoring is broken.

\paragraph{Future ablation models.}
\(M_0\) is the tied non-MoE control.  \(M_{\rm MoE}\) uses a plain softmax or sparse top-\(r\)
router and marginal-load balance only for S+1; \(M_{\rm static}\) freezes routes across depth.
\(M_{\rm nonrev}\) stores activations and isolates the reversible-memory benefit.  These are not
part of the current P1 executable interface.  Their code path lives in
\code{experiments/p1\_feedback\_stages.py} and is diagnostic until the positive control and \(M_0\)
make depth utility detectable.

\paragraph{Explicit exclusions from P1.}
The following are not part of the core science path: Dirichlet--UCB router, int6 quantization, DEQ or
implicit backward, fixed-point consistency, prefix-anchor consistency, spectral or Lyapunov pressure,
UFID, distillation, and anytime losses.  They may remain in the repository as legacy, deployment, or
later-stage axes, but they should not be used to establish P1.  MLA, smooth-sparse MoE, and the MoS
output head are \emph{retained} in \(M_0\) for the resource and expressiveness goals
(\S\ref{sec:goal}, \S\ref{sec:related}); the MoE router type is a control-experiment knob, not an
excluded mechanism.

%==============================================================================
\section{Experimental Protocol}
\label{sec:protocol}

\paragraph{Tier 1: detectability on synthetic depth-hard tasks.}
The primary task is iterated \(S_5\) permutation composition, with sequence length \(\ell\) tuning
the required depth.  Barrington's theorem shows that bounded-width polynomial-size branching
programs recognize exactly nonuniform \(\mathrm{NC}^1\), and the \(S_5\) permutation construction
is the canonical reason this task is a principled depth probe \cite{barrington}.  It is not a
language-modeling benchmark; it is the depth-demand smoke test.  Alternates are parity as a sanity
floor and \(s\)-\(t\) graph reachability with tunable diameter.

Run \(M_0\), the vanilla positive control, and \(M_{\rm clk}\) through the identical harness.  The
purpose is to confirm that \(G_T>0\) is detectable and to localize failure before spending compute
on small language modeling.  The in-repository executable harness is
\code{experiments/p1\_synthetic.py}; its P1-promotion variants are \code{control}, \code{m0}, and
\code{mclk}.  The same runner deliberately rejects MoE, cache-slope, and int6 switches.  That
small interface is part of the experiment: a failed positive control should trigger task, budget,
or block-capacity debugging, not another mechanism.  The reproducible DDP pipeline is
\code{experiments/run\_p1\_synthetic\_pipeline.sh}.  The separate feedback-stage runner
\code{experiments/run\_feedback\_stage\_pipeline.sh} composes the S+0 run with later-stage
diagnostics; its outputs are not P1-promotion evidence.

\paragraph{Tier 2: deployment relevance.}
After Tier 1 passes, evaluate small language modeling on byte/char-level \code{text8} or
\code{enwik8}, or on TinyStories.  Report \bpb{} where natural, and compare against a matched small
dense transformer on the Pareto frontier.

\paragraph{Curriculum.}
Use Huginn-style recurrence-count sampling for \(q_{\rm deep}\) and test-time depth scaling, plus
per-step input injection \cite{huginn}.  Do not borrow truncated BPTT for the core: reversibility
affords full BPTT and is the memory advantage being tested.

\paragraph{Depth-gain measurement and the Occam anti-collapse fix.}
A tied recurrence can collapse to a fixed point---effective depth \(\varphi\to 0\), depth-gain
\(G_T\to 0\)---so the \code{M0} trainer first \emph{measures} whether depth buys anything before
trusting internal depth metrics.  The \code{--k-eval-sweep} flag evaluates one trained model at each
recurrent depth \(K\) after the final validation (rank-0 prints; all ranks run the loop because the
validation all-reduces) and emits \(G_T=\bpb(K_{\min})-\bpb(K_{\max})\) (positive when deeper
recurrence helps) plus the cheap eval-depth \(\varphi\) \emph{proxy} \code{phi\_eval} fitted over the
\(\{K:\text{loss}(K)\}\) map.  \code{phi\_eval} is an OLS log-slope of one trained model's loss versus
\emph{inference} depth \(K\) (test-time depth utilization) and is \emph{not} the principled Iso-Depth
recurrence-equivalence exponent; the latter is the train-\(r\) \(\varphi\) (\code{phi\_isodepth}) of
Eq.~\eqref{eq:isodepth}, fitted over models pretrained at different recurrence budgets \(r\) by
\code{experiments/measure\_phi.py} (\S\ref{sec:related}).
The sweep is computed from a \emph{single shared trajectory}: the recurrence is deterministic given
the seed \((a_0,b_0,x_0)\), so depth \(K_{\max}\) is depth \(K_{\min}\) \emph{continued}.  Rather than
re-running the recurrence from the seed for each \(K\) (cost \(\sum_K K\)), the trainer takes ONE
depth-\(K_{\max}\) pass and reads the reversible midpoint \(\tfrac12(a+b)\) off that trajectory at each
requested \(K\) (cost \(\approx K_{\max}\), roughly half for a geometric ladder such as
\(\{8,16,32,64\}\)).  This is also the \emph{cleanest paired depth comparison}: every \(K\) shares one
trajectory and therefore one recurrence-seed noise draw, so depth is inherently the only variable
(\(K_{\max}=K_{\min}\) continued) with no per-\(K\) RNG reset needed in the sweep.  The evaluation
batches are still materialized once and replayed (the streaming loader would otherwise score each
\(K\) on different samples).  The depth-gain is thus a reproducible controlled measurement on fixed
weights, and the per-pass RNG reset is retained only inside the finite-horizon shallow/deep hinge.
The \(\varphi\) fit drops non-finite depths so a single
diverged \(K\) cannot \(\mathrm{NaN}\)-poison the proxy.  When collapse is detected, the
simplest principled fix---tested first per Occam---is Huginn-style \emph{random state
initialization}~\cite{huginn} (\code{--init-state random}): the recurrence is seeded with small
random non-learnable \(a_0,b_0\) independent of the injected token state \(x_0\) (the default
\code{x0} keeps \(a_0=b_0=x_0\), byte-identical to the prior behavior), while \(x_0\) is still
injected at every step.  Combined with the input injection and depth sampling already present, a
path-independent seed forces the \(K\) steps to do real work mapping noise to the solution, so the
recurrence \emph{uses} depth rather than settling at a fixed point.  Both the exact fp64
reconstruction and the grad-equivalence of the reversible backward to ordinary BPTT are preserved
for both seed modes: the algebraic inverse recovers the random seed exactly, and the custom backward
drops the \(x_0\)-seed gradient term (\(g_{x_0}\mathrel{+}=g_a+g_b\)) when the seed is the random,
non-learnable init---returning \(g_a,g_b\) in the seed slots, which the caller discards.

\paragraph{The principled expressiveness metric: depth-resolved \texorpdfstring{$\mathcal{V}$}{V}-usable predictive information.}
\label{par:vinfo}
The headline measure of the \emph{expressiveness} a model \emph{gains from recurrence depth} \(K\) is
the depth-resolved \(\mathcal{V}\)-usable predictive information~\cite{vinfo,vusable}.  Let \(Y\) be
the next token and \(z_K\) the depth-\(K\) representation read by the model's own readout family
\(\mathcal{V}\) (the MoS head).  Define the predictive \(\mathcal{V}\)-information
\begin{equation}
\label{eq:vinfo}
  I_{\mathcal{V}}(z_K\!\to\!Y)
  \;=\; H_{\mathcal{V}}(Y)\;-\;H_{\mathcal{V}}(Y\mid z_K)
  \;=\; H_{\mathcal{V}}(Y)\;-\;\mathrm{loss}(K),
\end{equation}
where \(H_{\mathcal{V}}(Y)=-\sum_y p(y)\log p(y)\) is the best \emph{marginal} predictor loss
(the empirical next-token marginal cross-entropy over the eval set, a constant in \(K\); reported in
nats and, dividing by \(\ln 2\), in bits), and \(\mathrm{loss}(K)\) is the depth-\(K\) conditional
cross-entropy already produced by the \code{--k-eval-sweep}.  The objective is
\begin{equation}
\label{eq:vinfo-obj}
  \mathrm{corr}\bigl(I_{\mathcal{V}}(K),\,K\bigr)\;\to\;+1
  \qquad\Longleftrightarrow\qquad \mathrm{loss}(K)\ \text{monotonically decreasing in }K,
\end{equation}
which the trainer logs as the rank correlation \code{expressiveness\_rho}
\(=\mathrm{spearman}\bigl(K,\,I_{\mathcal{V}}(K)\bigr)\) (equivalently of \(-\mathrm{loss}\) vs.\ \(K\)),
together with the per-\(K\) \code{usable\_info} line (\(I_{\mathcal{V}}(K)\) in bits) and the total
usable information gained across the swept range \code{iv\_total\_bits}
\(=I_{\mathcal{V}}(K_{\max})-I_{\mathcal{V}}(K_{\min})\) (in bits).  On real data \(I_{\mathcal{V}}(K)\)
is \emph{realized} expressiveness; on the depth-hard positive control it measures \emph{capability}.

This is THE principled expressiveness measure for three reasons.  (i) \(\mathcal{V}\)-information is
the \emph{unique} information measure that can be \emph{created by computation}: Shannon mutual
information cannot grow with depth (the data-processing inequality forbids
\(I(z_K;Y)>I(z_0;Y)\) for any deterministic map), but \(I_{\mathcal{V}}\) can, because it is defined
relative to the bounded readout family \(\mathcal{V}\)~\cite{vinfo}.  (ii) It is
\emph{sufficient-and-necessary} for usable predictive expressiveness w.r.t.\ \(\mathcal{V}\): it equals
the codelength (MDL) reduction a \(\mathcal{V}\)-decoder achieves from \(z_K\), and captures the
combinatorial MoE path family (a parameter-count surrogate \(\varphi N\) cannot).  (iii) It
\emph{subsumes} the previous depth summaries: the depth-gain \(G_T=\bpb(K_{\min})-\bpb(K_{\max})\) is a
two-point bits-per-byte version of \code{iv\_total\_bits}, and \code{phi\_eval} is an OLS log-slope of
the same \(\{K:\mathrm{loss}(K)\}\) curve --- both are now read as summaries/proxies of
Eq.~\eqref{eq:vinfo}.  The marginal entropy \(H_{\mathcal{V}}(Y)\) is accumulated once as token counts
over the same cached eval batches and all-reduced across DDP ranks so the constant is identical on
every rank; the sweep drops non-finite depths, and \code{expressiveness\_rho} is \(\mathrm{NaN}\) with
fewer than three finite \(K\) (no rank correlation estimable).

\paragraph{Monotonicity guarantee and its trainable surrogate.}
Optimizing \(I_{\mathcal V}\) provably yields a more task-predictive model --- it \emph{is} the achievable
readout-loss reduction --- and reversibility guarantees no input information is \emph{destroyed} across
depth (\(z_K\) is a bijection of \((z_0,x_0)\); \code{recon\_rel}\(\approx0\)).  A \emph{monotone}
guarantee, \(I_{\mathcal V}(K)\ge I_{\mathcal V}(K{-}1)\) (i.e.\ \code{expressiveness\_rho}\(=+1\) by
construction), additionally requires a \emph{nested} achievable class
\(\mathcal F_{K-1}\subseteq\mathcal F_{K}\) so that depth \(K\) can always reproduce depth \(K{-}1\).
Three structures supply this: an \emph{anytime} (depth-selecting / nested) readout supplies it directly;
a \emph{convergent} contraction (a Huginn fixed point) supplies it asymptotically; \emph{untied} layers
supply it by class inclusion.  The current \emph{tied}, finite-horizon readout \(z_K\to\textsc{MoS}\)
does \emph{not} nest by construction --- with weights tied across steps one cannot make the \(K\)-th step
the identity without making \emph{all} steps the identity --- so deep \(K\) can over-think (degrade)
beyond the trained horizon.  We therefore enforce monotonicity \emph{softly}, via the stop-gradient
no-degradation hinge \(\lambda_h\,[\,L_T(K_{\rm hi})-\sg\,L_T(K_{\rm lo})+m\,]_+\) plus Huginn-style depth
sampling, and \emph{verify} it empirically through \code{expressiveness\_rho}; on a depth-hard positive
control the circuit-complexity equivalence (a bounded model solves a depth-\(d\) serial task iff it
realizes serial computation depth \(\ge d\))~\cite{cot} makes the increase \emph{strict}.  A hard
guarantee remains an available design lever (anytime readout, convergence, or untying), each trading
against the finite-horizon / weight-tying resource win.

\paragraph{Expressiveness versus efficiency: orthogonal goals on the two-goal frontier.}
The recurrence param-\emph{efficiency} exponent \(\varphi_{\rm isodepth}\) (Eq.~\eqref{eq:isodepth}) and
\(\mathcal{V}\)-information \(I_{\mathcal{V}}\) (Eq.~\eqref{eq:vinfo}) measure \emph{orthogonal} things.
\(\varphi_{\rm isodepth}\) is a \emph{resource-axis} quantity --- how many \emph{unique parameters} one
recurrent loop is worth (\(\varphi=1\): a loop equals a unique block; \(\varphi>1\):
\emph{super}-param-equivalent recurrence, e.g.\ via MoE combinatorial path reuse, so the fit bound is
relaxed to \(\varphi\in[0,3]\) and no longer censored at unity).  \(I_{\mathcal{V}}\) is a
\emph{quality-axis} quantity --- how much \emph{usable predictive information} depth actually adds,
independent of parameter accounting.  A model can be param-efficient yet low-expressiveness (high
\(\varphi\), flat \(I_{\mathcal{V}}\)) or expressive yet param-inefficient.  Both are gated by the same
reversibility correctness condition \(\code{\detokenize{recon_rel}}\approx0\): the depth-\(K\)
representation \(z_K\) and the loop param count are only meaningful when the reversible round-trip is
bit-exact.

\paragraph{Scale and statistics.}
Use \(d_{\rm model}\in[512,768]\) (the trainer default is \(768\): 12 query / 4 KV heads,
head dim 64, MLA kv-latent 256, 16 experts of rank 64), one reversible block, context length at
most 1024, and multiple seeds.  Evaluate
\[
  K\in\{4,8,16,32,64,128\}.
\]
Report the convergence knee first, then choose \((K_{\rm lo},K_{\rm hi})\) to straddle it.  Every
headline \(G_T\) uses paired confidence intervals.  The executable DDP harness emits a streaming
paired normal-approximation interval from all-reduced sums and squared sums; offline reports may
replace this with a paired bootstrap over saved per-example gains.

%==============================================================================
\section{Metrics: Gates Versus Diagnostics}
\label{sec:metrics}

This section is the authoritative reference for every quantity the trainer emits.  Each metric is
classified along two axes.  \emph{Role}: a \GATE\ decides promotion (a hard pass/fail), whereas a
\DIAG{} explains \emph{why} a gate passed or failed but never promotes on its own.  \emph{Goal axis}:
the project optimizes two \emph{orthogonal} objectives --- \textbf{expressiveness} (how much usable
predictive information depth adds, maximize) and \textbf{resource} (activation memory, minimize / hold
flat in \(K\)) --- bridged by an \textbf{efficiency-vs-dense} axis (the Iso-Depth exponent \(\varphi\))
and underwritten by a single \textbf{correctness} gate (\code{\detokenize{recon_rel}}) that makes
every other number trustworthy.  The reader should read each metric's mathematical definition before
its logged value (the meta-principle: definition first, value second).

\subsection{Summary: which metrics gate, which diagnose}
\renewcommand{\arraystretch}{1.2}
\begin{center}
\footnotesize
\begin{tabularx}{\textwidth}{@{}L{0.27\textwidth} c L{0.50\textwidth}@{}}
\toprule
\textbf{Metric} & \textbf{Role} & \textbf{Axis / one-line rule} \\
\midrule
\(\code{\detokenize{recon_rel}}\) & \GATE\ (correctness) & Reversible round-trip error \(<10^{-6}\); makes all other metrics trustworthy (\S\ref{ssec:correctness}). \\
\(G_T(K_{\rm lo},K_{\rm hi})\), \(I_{\mathcal V}\), \(\code{\detokenize{expressiveness_rho}}\) & \GATE\ (P1, expressiveness) & Paired CI lower bound \(>0\) on a depth-hard task, with knee-straddling pairs and a passing positive control; \(I_{\mathcal V}\) is the principled form, \(G_T\) its two-point bpb summary (\S\ref{ssec:expressiveness}). \\
\(\mathrm{NDR}_\epsilon\) & \GATE\ (guard) & No-degradation rate \(\le\delta_{\rm task}\) (depth must not \emph{hurt}). \\
\(\vram_{\rm peak}\), \(\mathrm{slope}(\vram/\text{batch})\), \(R_{\rm act}(K)\) & \GATE\ (P2, resource = memory) & Peak VRAM low \emph{and flat in \(K\)}; \(R_{\rm act}(K)\le1.10\) (\S\ref{ssec:resource}). \\
Pareto \((S_T,\vram,\mathrm{ms/token})\) & \GATE\ (sufficiency) & On or above frontier versus a matched small dense transformer.  This, not \(G_T>0\) alone, is project success. \\
\midrule
\(\varphi_{\rm isodepth}\) (Recurrence-equivalence exponent) & \DIAG\ (efficiency-vs-dense, orthogonal) & How many \emph{unique-parameter} blocks \(r\) loops are worth; \(\varphi\in[0,3]\) (\(>1\): super-param-equivalent).  Resource/efficiency axis, \emph{not} expressiveness (\S\ref{ssec:efficiency}). \\
\(\code{\detokenize{phi_eval}}\) & \DIAG\ (efficiency proxy) & Cheap single-model eval-depth log-slope proxy; \emph{not} \(\varphi_{\rm isodepth}\) (\S\ref{ssec:efficiency}). \\
\(\code{\detokenize{erank}}\), \(\code{\detokenize{disp_tail}}\) & \DIAG\ (expressiveness) & Representation dimensionality (non-monotone) and sustained-dynamics / effective-depth signals (\S\ref{ssec:expressiveness}). \\
\(\code{\detokenize{route_step_div}}\), \(\code{\detokenize{expert_cos_div}}\), \(\code{\detokenize{router_entropy}}\), \(\code{\detokenize{expert_util}}\) & \DIAG\ (MoE-basis engagement) & Across-step path diversity, basis non-degeneracy, specialization, utilization (\S\ref{ssec:moe}). \\
\(\code{\detokenize{kv_bytes}}\), \(\code{\detokenize{ops:vram_util_pct}}\), \(\code{\detokenize{ops:tok_per_s}}\) & \DIAG\ (resource / ops) & MLA KV footprint; VRAM-util and throughput tuning knobs --- explicitly \emph{not} gates (\S\ref{ssec:resource}). \\
Convergence-depth \(k^\star(x)\), \(\mathrm{WFIR}\), \(\Delta_{\rm MoE}\), \(\mathrm{AEBR}\), route-depth NMI, \(\rho_F\)/\(\sigma_{\max}(J_F)\), \(\mathrm{Gap}_{\rm dense}\) & \DIAG\ (legacy / future P3) & Failure-analysis traces; not emitted by the P1 harness and not P1 gates (\S\ref{ssec:demoted}). \\
\bottomrule
\end{tabularx}
\end{center}

\subsection{Correctness gate}
\label{ssec:correctness}
\paragraph{\(\code{\detokenize{recon_rel}}\) --- reversible round-trip error (\GATE).}
Let \(\mathrm{forward}(\cdot)\) run the depth-\(K\) recurrence from a seed \((a_0,b_0)\) to the terminal
\((a_K,b_K)\) and \(\mathrm{invert}(\cdot)\) apply the algebraic inverse of
\S\ref{sec:models}, producing \((a_0',b_0')=\mathrm{invert}(\mathrm{forward}(a_0,b_0))\).  Then
\begin{equation}
\label{eq:recon-rel}
  \code{\detokenize{recon_rel}}
  \;=\;
  \frac{\lVert a_0'-a_0\rVert+\lVert b_0'-b_0\rVert}{\lVert a_0\rVert+\lVert b_0\rVert+\epsilon}.
\end{equation}
\emph{Interpretation.}  The \(O(1)\)-in-\(K\) reversible backward (\code{RevRecurrenceFn}) does not
store intermediate activations; it \emph{reconstructs} them from \((a_K,b_K)\) via exactly this inverse,
so the reconstructed graph --- and hence the BPTT gradients --- equal ordinary stored-activation BPTT
\emph{iff} the round-trip is exact.  A large \(\code{\detokenize{recon_rel}}\) means the reconstructed
graph (and the gradient signal) has drifted from the true forward.  It is the gate that makes every
other metric trustworthy: \(z_K\), the loss-vs-\(K\) curve, and the parameter count are only meaningful
when the gradients that produced them were correct.  Measured at the rank-0 log site \emph{under the
training bf16 autocast} (one forward + one inverse sweep), so it reports actual training reconstruction
drift; fp64 stream accumulation (\S\ref{sec:models}, \code{\detokenize{recurrence_accum_dtype}}) and
autocast replay drive it to \(\approx0\).  \emph{Goal:} \(<10^{-6}\) (machine precision in fp64).
\GATE\ (correctness; also a P2 sub-gate).

\subsection{Expressiveness axis (maximize)}
\label{ssec:expressiveness}
The headline goal is that recurrent depth \(K\) adds \emph{usable} predictive information.

\paragraph{\(I_{\mathcal V}(z_K\!\to\!Y)\) --- \(\mathcal V\)-usable predictive information (\GATE\ basis).}
With \(Y\) the next token and \(z_K\) the depth-\(K\) representation read by the model's own readout
family \(\mathcal V\) (the MoS head; \S\ref{sec:models}),
\begin{equation}
\label{eq:iv-metric}
  I_{\mathcal V}(z_K\!\to\!Y)
  \;=\;
  H_{\mathcal V}(Y)-H_{\mathcal V}(Y\mid z_K)
  \;=\;
  H_{\mathcal V}(Y)-\mathrm{nll}(z_K),
  \qquad
  H_{\mathcal V}(Y)=-\sum_y\hat p(y)\log\hat p(y),
\end{equation}
where \(H_{\mathcal V}(Y)\) is the empirical next-token \emph{marginal} cross-entropy (the best
marginal predictor's loss, a constant in \(K\), accumulated once over the cached eval targets and
all-reduced across DDP ranks; \(\hat p\) the empirical token frequency) and \(\mathrm{nll}(z_K)\) is the
depth-\(K\) conditional cross-entropy from the \code{--k-eval-sweep}.  Emitted in bits as
\(\code{\detokenize{iv_bits}}=I_{\mathcal V}/\ln 2\) per \(K\) (the \code{\detokenize{usable_info}}
line).  \emph{Why it is the principled depth-expressiveness measure} (\cite{vinfo,vusable}): \(\mathcal
V\)-information can be \emph{created by computation} and so can \emph{grow with} \(K\), unlike Shannon
mutual information, which the data-processing inequality forbids from increasing under the deterministic
map \(z_0\!\to\!z_K\) (\(I(z_K;Y)\le I(z_0;Y)\)); it is defined relative to the bounded readout family
\(\mathcal V\), making it sufficient-and-necessary for usable predictive expressiveness w.r.t.\ \(\mathcal
V\), and it equals the MDL codelength reduction a \(\mathcal V\)-decoder achieves from \(z_K\).
\emph{Goal:} \(I_{\mathcal V}(K)\) increasing in \(K\).

\paragraph{\(\code{\detokenize{expressiveness_rho}}\) --- headline expressiveness gate-direction (\GATE).}
\begin{equation}
\label{eq:expr-rho}
  \code{\detokenize{expressiveness_rho}}
  \;=\;
  \rho_{\rm Spearman}\bigl(K,\,I_{\mathcal V}(K)\bigr)
  \;=\;
  \rho_{\rm Spearman}\bigl(K,\,-\mathrm{nll}(K)\bigr)\in[-1,1].
\end{equation}
\(=+1\) iff depth monotonically adds usable information (the objective); \(\mathrm{NaN}\) with fewer than
three finite \(K\) (no rank correlation estimable).  This is the headline expressiveness reading;
gate-direction \(\to+1\).

\paragraph{\(\code{\detokenize{iv_total_bits}}\) --- total depth-added usable information (\DIAG).}
\(\code{\detokenize{iv_total_bits}}=I_{\mathcal V}(K_{\max})-I_{\mathcal V}(K_{\min})\) bits: the net
usable information depth buys across the swept range.

\paragraph{\(G_T\) --- depth-gain (\GATE\ summary).}
\(G_T(K_{\rm lo},K_{\rm hi})=\bpb(K_{\min})-\bpb(K_{\max})\) (bits-per-byte units), the two-point bpb
summary of \(I_{\mathcal V}\) and the operational P1 statistic.  \emph{Gate rule:} paired-CI lower
bound \(>0\) on a depth-hard task with knee-straddling \((K_{\rm lo},K_{\rm hi})\) and a passing
positive control; the executable harness emits a streaming paired normal-approximation interval from
all-reduced sums and squared sums.  \GATE\ (P1).

\paragraph{\(\mathrm{NDR}_\epsilon\) --- no-degradation rate (\GATE\ guard).}
The fraction of paired examples on which deeper recurrence \emph{degrades} the score by more than a
task tolerance \(\epsilon\); a guard that depth must not hurt, with rule \(\mathrm{NDR}_\epsilon\le
\delta_{\rm task}\).

\paragraph{\(\code{\detokenize{erank}}\) --- effective (numerical) rank (\DIAG).}
For singular values \(\sigma_i\) of a probed matrix (e.g.\ the tied embedding / expert output space),
with \(\hat\sigma_i=\sigma_i/\sum_j\sigma_j\),
\begin{equation}
\label{eq:erank}
  \code{\detokenize{erank}}
  \;=\;
  \exp\!\Bigl(-\textstyle\sum_i\hat\sigma_i\log\hat\sigma_i\Bigr)
\end{equation}
(Roy--Vetterli~\cite{effrank}; the exponential of the spectral entropy, \(=n\) for an isotropic
spectrum, \(=1\) for rank-1).  A representation-dimensionality \emph{diagnostic}.  It is \emph{explicitly
non-monotone} and \emph{not} a maximize-target: compression of the representation often aids
generalization (the intrinsic-dimensionality caveat), so a falling \code{\detokenize{erank}} is not by
itself a regression.

\paragraph{\(\code{\detokenize{disp_tail}}\) --- sustained-dynamics / effective-depth (\DIAG).}
With \(z_k=\tfrac12(a_k+b_k)\), the normalized per-step displacement is \(\lVert z_{k+1}-z_k\rVert/\lVert
z_k\rVert\); \code{\detokenize{disp_tail}} is its tail mean over the last half of the budget,
\begin{equation}
\label{eq:disp-tail}
  \code{\detokenize{disp_tail}}
  \;=\;
  \frac{2}{K}\!\!\sum_{k>K/2}\frac{\lVert z_{k+1}-z_k\rVert}{\lVert z_k\rVert}.
\end{equation}
A sustained (non-decaying) tail signals the recurrence is still doing work deep in the budget (high
effective depth); a decay to \(\approx0\) signals fixed-point saturation (the collapse mode the
\code{--init-state random} and \code{--damping orthogonal} fixes target).

\subsection{Efficiency-vs-dense axis (the \texorpdfstring{$\varphi$}{phi} bridge)}
\label{ssec:efficiency}
This axis is \emph{complementary} to \(I_{\mathcal V}\): it measures how \emph{parameter-efficient} the
recurrence is relative to a dense deep model, independent of absolute usable information.

\paragraph{\(\varphi_{\rm isodepth}\) --- the Iso-Depth recurrence-equivalence exponent (\DIAG, the principled \(\varphi\)).}
This is \emph{the} \(\varphi\), fitted directly from the Iso-Depth joint scaling law~\cite{isodepth}
(\S\ref{sec:related}, Eq.~\eqref{eq:isodepth}):
\(L(r)=E+A\,(N_{\rm once}+r^{\varphi}N_{\rm rec})^{-\alpha}\), over a recurrence-count sweep
\(r\in\{1,2,4,8,16\}\) of \emph{looped models} (\(N_{\rm once}\): non-recurrent params; \(N_{\rm
rec}\): recurrent/looped-block params), with \(\{E,A,\alpha,\varphi\}\) fitted by nonlinear least
squares (\(\varphi\in[0,3]\)).  \emph{Interpretation:} \(r\) tied loops \(\equiv r^{\varphi}\) unique
blocks --- \(\varphi=1\) full equivalence (tying is free), \(\varphi<1\) tying costs capacity (the paper
measured \(\varphi\approx0.46\)), \(\varphi>1\) \emph{super}-param-equivalent (the MoE-basis
``exponential paths'' beating dense depth).  This is the efficiency-vs-dense axis (a resource-axis
quantity), \emph{not} expressiveness.  Measured by \code{experiments/measure\_phi.py} as a multi-\(r\)
\emph{training} sweep (an experiment-level estimate, not a per-run log line); the five-point fit is
small-sample and noise-fragile, so a trustworthy estimate needs more \(r\) and/or repeated seeds.

\paragraph{\(\code{\detokenize{phi_eval}}\) --- cheap single-model eval-depth proxy (\DIAG, \emph{not} \(\varphi_{\rm isodepth}\)).}
A per-run proxy fitted from \emph{one} trained model's \code{--k-eval-sweep}: model
\(\mathrm{nll}(K)\approx a-b\log K\) and set \(\varphi_{\rm eval}=\mathrm{clamp}(b/b_{\rm ref},0,1)\)
(\(b_{\rm ref}=1\)).  It is an OLS log-slope of loss versus \emph{inference} depth \(K\) (test-time depth
utilization) with \emph{no} train-\(r\) sweep and \emph{no} \((N_{\rm once},N_{\rm rec})\) split, so it
is explicitly a \emph{proxy} --- useful for a quick per-run read, but distinct from and not a substitute
for the principled \(\varphi_{\rm isodepth}\) of Eq.~\eqref{eq:isodepth}.

\subsection{MoE-basis engagement (diagnostics)}
\label{ssec:moe}
These probe whether the smooth-sparse MoE forms a diverse depth-\(K\) path family (the P3 mechanism);
all are log-site reads (no hot-loop sync).

\paragraph{\(\code{\detokenize{route_step_div}}\) --- across-step expert diversity.}
For each token, count the distinct dominant experts \(\arg\max_e p^{(k)}_e\) across the \(K\) recurrence
steps, divided by \(K\), then mean over tokens:
\begin{equation}
\label{eq:route-step-div}
  \code{\detokenize{route_step_div}}
  \;=\;
  \E_t\!\left[\frac{\#\{\,\arg\max_e p^{(k)}_{t,e}\;:\;k=1,\dots,K\,\}}{K}\right]\in[1/K,\,1].
\end{equation}
\(=1\): a different expert each step (the exponential-path basis is engaged); \(=1/K\): the same expert
every step (MoE-basis unused, the across-depth collapse mode).

\paragraph{\(\code{\detokenize{expert_cos_div}}\) --- basis non-degeneracy.}
Mean pairwise \(1-\cos\) of the experts' output vectors on a shared probe batch, \(\in[0,1]\); high
\(\Rightarrow\) diverse expert functions, \(\approx0\Rightarrow\) duplicated/degenerate experts.

\paragraph{\(\code{\detokenize{router_entropy}}\) --- specialization.}
Mean per-token router entropy \(H(p)=-\sum_e p_e\log p_e\); \emph{low} \(\Rightarrow\) peaked/specialized
routing.  This is specialization, \emph{not} utilization.

\paragraph{\(\code{\detokenize{expert_util}}\) --- utilization.}
Global entropy \(H(\bar p)\) of the batch-mean routing \(\bar p=\E_t\,p_t\), \(\in[0,\ln E]\); near
\(\ln E\) \(\Rightarrow\) all experts used, low \(\Rightarrow\) dead experts.  \(\code{\detokenize{diag:active_frac}}\)
(mean routed-weight active fraction) is a related MoE-mechanism diagnostic, \emph{not} a resource gate.

\subsection{Resource axis (minimize / hold flat in \texorpdfstring{$K$}{K})}
\label{ssec:resource}
The resource goal is \emph{activation memory}, which the reversible \(O(1)\)-in-\(K\) backward holds
flat in recurrent depth.  Dropped as resource metrics: the artifact-16\,MB hard gate (artifact bytes
logged informationally only), FLOPs/token, and wall-clock/step.

\paragraph{\(\code{\detokenize{peak_vram}}\) --- primary resource number (\GATE).}
Peak allocated VRAM in MiB (\(\max\) over the step).  \emph{Goal:} low \emph{and flat in \(K\)} --- the
reversible recurrence's defining property.  \GATE\ (P2).

\paragraph{\(R_{\rm act}(K)=\mathrm{mem}(2K)/\mathrm{mem}(K)\) --- activation-memory ratio (\GATE).}
The ratio of activation memory at doubled depth; \(R_{\rm act}(K)\le1.10\) certifies activations are
\(\approx\)constant in recurrent depth, the P2 claim.  Supplied by the control sweep.  \GATE\ (P2).

\paragraph{VRAM-vs-batch slope --- \(\mathrm{slope}(\vram/\text{batch})\) (\GATE).}
The OLS slope (MB/sample) of peak VRAM versus batch size (\code{\detokenize{vram_vs_batch_scaling}});
goal low/flat.  \code{\detokenize{--grad-accum 1}} runs the whole batch in one forward to fill VRAM,
which the \(O(1)\)-in-\(K\) reversible recurrence makes affordable at deep \(K\).

\paragraph{\(\code{\detokenize{kv_bytes}}\) --- MLA KV-cache footprint (\DIAG).}
Bytes per token of the MLA-compressed KV cache (the latent plus decoupled RoPE channel, billed bf16);
the static deployment footprint.  The autoregressive constant-KV claim (\(\alpha_{\rm kv}\approx0\)) is
an S+2 axis (\S\ref{sec:staging}), not a current gate.

\paragraph{\(\code{\detokenize{ops:vram_util_pct}}\), \(\code{\detokenize{ops:tok_per_s}}\) --- ops/efficiency tuning (\DIAG, \emph{not} gates).}
VRAM utilization \(\code{\detokenize{ops:vram_util_pct}}=100\cdot\vram_{\rm peak}/\)device memory and
throughput \(\code{\detokenize{ops:tok_per_s}}=\text{batch tokens}/\text{step wall-time}\), logged under
the \code{ops:} prefix for tuning the batch / grad-accum knobs.  They are \emph{explicitly not}
resource gates.

\subsection{Two-goal frame}
\label{ssec:two-goal}
Expressiveness (\(I_{\mathcal V}\) / \code{\detokenize{expressiveness_rho}}, \S\ref{ssec:expressiveness})
and resource (\(\code{\detokenize{peak_vram}}\) / \(R_{\rm act}\), \S\ref{ssec:resource}) are
\emph{orthogonal} goals on the Pareto frontier: a model can be expressive yet memory-hungry, or
memory-flat yet flat in usable information.  \(\varphi_{\rm isodepth}\) (\S\ref{ssec:efficiency}) is the
efficiency-vs-dense bridge between them (how few unique parameters one loop is worth), and the
correctness gate \(\code{\detokenize{recon_rel}}\approx0\) (\S\ref{ssec:correctness}) underwrites
\emph{both} axes --- \(z_K\) and the loop param count are only meaningful when the reversible round-trip
is bit-exact.  Project success is the Pareto gate, not \(G_T>0\) alone.

\subsection{Demoted / legacy internal metrics (diagnostics, not P1 gates)}
\label{ssec:demoted}
The following are failure-analysis traces, not P1 gates, and are not emitted by the P1 harness.

\paragraph{Useful function-space depth (\(\mathrm{WFIR}\)).}
For a fixed probe set \(X\), let \(f_k(X)\) be a task-aligned readout at recurrent depth \(k\), with
increments \(\Delta_k(X)=f_{k+1}(X)-f_k(X)\).  \(\mathrm{WFIR}\) is the effective rank of centered,
regularized functional increments, reported against a shuffled-depth baseline
(\(\mathrm{WFIR}_{\rm real}-\mathrm{WFIR}_{\rm shuffled}\)); use a task-aligned readout (correct-token
log-probability or NLL contribution), never random projections, as the headline.  It is diagnostic
because diversity alone can be noise; it becomes meaningful only when paired with positive \(G_T\) and
controlled \(\mathrm{NDR}_\epsilon\).

\paragraph{Other legacy traces.}
Convergence-depth \(k^\star(x)\) versus task hardness (flat-and-small \(k^\star\) is the plateau
signature); the P3 MoE-ablation metrics \(\Delta_{\rm MoE}\), \(\mathrm{AEBR}\), and route-depth NMI
(used only after P1 passes, in a separate MoE ablation, never emitted by the P1 harness); the
fixed-point convergence diagnostics \(\rho_F\), residual, and \(\sigma_{\max}(J_F)\) (advisory, feeding
\(k^\star\) and fallback decisions); and the dense-simulation diagnostics \(\mathrm{Gap}_{\rm dense}\),
\(\mathrm{KL}_{\rm teacher}\) (the weakest claim, never a gate).  The older internal effective-depth
traces \code{\detokenize{ED_update}}, \code{\detokenize{ED_logit}},
\code{\detokenize{route_depth_nmi_mean}}, \code{\detokenize{route_depth_nmi_max}}, and
\code{\detokenize{expert_util_mean}} map to ED\(_u\), ED\(_\ell\), route-depth NMI, and
expert-utilization traces.  They are useful for failure analysis, but they are \emph{not P1 gates}.

%==============================================================================
\section{Clock Diagnostic and Failure Triage}
\label{sec:triage}

Run \(M_0\), the vanilla control, and \(M_{\rm clk}\) on Tier 1.  Interpret outcomes as follows:
\renewcommand{\arraystretch}{1.2}
\begin{center}
\footnotesize
\begin{tabularx}{\textwidth}{@{}L{0.46\textwidth} L{0.46\textwidth}@{}}
\toprule
\textbf{Observation} & \textbf{Localized cause / action} \\
\midrule
Control does not scale on a provably depth-hard task & Harness or metric is broken: pairs may not straddle the knee, scoring may be wrong, or the task is too easy.  Fix this before model work. \\
Control scales; \(M_0\) scales & Core is sound.  Proceed to Tier 2, then optional P3 work. \\
Control scales; \(M_0\) does not; \(M_{\rm clk}\) scales & Bottleneck is autonomous depth signal, such as vanishing displacement.  Pursue sustained dynamics, velocity-state routing, or accept a mild clock with measured extrapolation cost. \\
Control scales; neither \(M_0\) nor \(M_{\rm clk}\) scales & Bottleneck is block expressivity or parameterization, not the signal.  Repair the block, not the router. \\
\(M_{\rm clk}\) raises \(G_T\), but extrapolates poorly & Clock works as an upper bound and its cost is extrapolation.  This justifies investing in an autonomous substitute. \\
\bottomrule
\end{tabularx}
\end{center}
The clock is a diagnostic ceiling, not a design commitment.  It answers whether depth-conditioned
computation is useful on the data before investing in autonomous routing.

\paragraph{Autonomous substitute: a strictly-invertible damping (hyper-connection) matrix.}
When the clock \(M_{\rm clk}\) scales but \(M_0\) does not (autonomous-depth-signal bottleneck), the
principled escalation is \emph{not} to ship the clock---which a learned per-step embedding
\(e_k\) makes depth-conditioned and therefore poorly-extrapolating, the very property Huginn's
depth-unconditioned design avoids \cite{huginn}---but to give the recurrence \emph{autonomous}
sustained dynamics via a strictly-invertible mixing/``damping'' matrix on the coupling streams.
This is the reversible specialization of Hyper-Connections \cite{hyperconnections}, whose
unconstrained mixing matrix is highly expressive (parallel residual streams \(\to\) high effective
depth) but unstable (signal blow-up).  The stable form constrains the matrix to a manifold:
manifold-constrained Hyper-Connections \cite{mhc} use a doubly-stochastic (Birkhoff) matrix, and
the \emph{orthogonal} variant \cite{jpmhc} gives dynamical isometry.  For our reversible core the
orthogonal choice is canonical: with \(Q\) orthogonal the coupling
\(a_{k+1}=Q\,a_k+F(b_k{+}x_0),\; b_{k+1}=Q\,b_k+G(a_{k+1}{+}x_0)\) inverts \emph{exactly} as
\(Q^{\top}\) (no Parcae reverse-reconstruction floor), is norm-preserving (so per-step displacement
cannot decay to a fixed point and the dynamics stay well-scaled \emph{beyond} the trained depth,
unlike the clock), and the cross-stream/cross-depth mixing supplies expressiveness.  It thus
supersedes both the UT-style clock and the diagonal Parcae damping-with-floor in one mechanism,
relaxing the \(A{+}B{=}I\) convex tie of Parcae damping to an orthogonal \(A=Q\) with free
\(B\) (the residual updates \(F,G\)).
Per Occam it is the \emph{second} rung---tested only if Huginn-style random init
(\code{--init-state random}, \S\ref{sec:metrics}) does not by itself close the \(M_0\)/\(M_{\rm clk}\)
gap.

\emph{Implemented.}  This is now available behind the \code{--damping} knob (\code{none}, the
default, is \(Q{=}I\) and byte-identical to the additive coupling; \code{orthogonal} enables the
mechanism) with rank knob \code{\detokenize{--damping-rank}} (default \(r{=}16\)).  \(Q\) is
parameterized as the \emph{matrix exponential of a low-rank skew}, \(Q=\expm(S)\), of a
skew-symmetric \(S=UV^{\top}-VU^{\top}\) with learnable \(U,V\in\mathbb{R}^{d\times r}\);
\(S^{\top}{=}{-}S\) guarantees \(\expm(S)\) is \emph{exactly} orthogonal for \emph{any} \(\|S\|\)
(\(Q^{\top}Q{=}I\), \(\det Q{>}0\), a proper rotation), and the inverse is exactly
\(Q^{\top}=\expm(-S)\) --- no floor and, crucially, \emph{no} \((I+S)^{-1}\) solve.  This replaces
the earlier low-rank skew-Cayley form \(Q=(I-S)(I+S)^{-1}\): the Cayley form is fp64-exact only
while \(\|S\|\) is small, but as the learnable \(U,V\) grow its low-rank Woodbury capacitance matrix
becomes ill-conditioned, the solve loses precision, \(Q\) drifts off the orthogonal manifold, and
\(Q^{\top}(Q\,x)\neq x\) --- so \code{\detokenize{recon_rel}} explodes mid-training and corrupts the
reconstructed gradients.  \(\expm(S)\) has no solve and is exactly orthogonal regardless of \(\|S\|\),
so \code{\detokenize{recon_rel}}\(\approx0\) is stable across training.  It is computed on the
low-rank active subspace to stay \(O(B\,T\,d\,r)\): with \(W=[U\,|\,V]\) and an orthonormal basis
\(B\) of \(\operatorname{span}(W)\) (a thin QR, \(k{\le}2r\) columns), the projected skew
\(S_{\rm sub}=B^{\top}SB\in\mathbb{R}^{k\times k}\) is exponentiated by a small \(k{\times}k\)
\code{\detokenize{matrix_exp}} (\(R=\expm(S_{\rm sub})\), \(R^{\top}=\expm(-S_{\rm sub})\)), and
\(Q\,x=x+B\big((R-I)(B^{\top}x)\big)\) (identity off \(\operatorname{span}(B)\), where \(S{=}0\), so
this is exactly the full-space \(\expm(S)\)); the \(d{\times}d\) matrix is \emph{never} formed.  \(Q\)
is shared across steps and between the \(a/b\) streams.  \(U,V\) are initialized at
\(10^{-3}\!\cdot\!\mathcal{N}\) so \(S{\approx}0,\ Q{\approx}I\) at start (near-identity readout);
rotation is learned.  \(Q\) is a fixed matrix \emph{within} a forward+backward (a function of \(U,V\)
only, not of the evolving states), so each reverse undo needs one already-available state (no
leapfrog circularity); applying \(Q\) in the fp64 accumulation stream keeps \(Q^{\top}(Q\,x){=}x\) to
\(\sim\!10^{-15}\), so the reversible round-trip \code{\detokenize{recon_rel}}\(\approx0\) even under
bf16 autocast at deep \(K\).  The custom \(O(1)\)-memory backward applies \(Q^{\top}\) in its
reconstruction and \(Q\) in the rebuilt step (and \(\expm\) / QR are differentiable), so \(U,V\)
receive gradients identical to ordinary BPTT.

%==============================================================================
\section{Assumptions and Checks}
\label{sec:assumptions}

\renewcommand{\arraystretch}{1.2}
\begin{center}
\footnotesize
\begin{tabularx}{\textwidth}{@{}c L{0.42\textwidth} L{0.40\textwidth}@{}}
\toprule
& \textbf{Assumption} & \textbf{How it is checked, not assumed} \\
\midrule
A1 & The task requires depth & Provable depth-hard synthetic task plus positive control.  If neither control nor clock scales, the harness or block is the issue. \\
A2 & Reversibility holds numerically & \(R_{\rm rec}\) measured with real stochasticity; replay RNG state for stochastic operations in the recurrent path. \\
A3 & Memory is binding and latency tolerable & Pareto gate versus a small dense model.  If a small dense model wins, the memory-bound niche is absent. \\
A4 & Nestedness holds without a halting gate & Depth is selected by the recurrent budget \(K\); verify the depth-\(K_{\rm lo}\) function is reachable inside the depth-\(K_{\rm hi}\) class, otherwise nestedness is only approximate. \\
A5 & Additive coupling stays bounded over \(K\) & Boundedness penalty plus hard guards; monitor norm growth as a stability and throughput cost. \\
A6 & Optional P3 capacity recovery is needed & Tested only after P1 holds; otherwise capacity mechanisms are premature. \\
\bottomrule
\end{tabularx}
\end{center}

The following are demoted from assumptions to tested hypotheses: routes specialize without a clock,
and dense soft routing is affordable.  The latter is moot in the P1 core and re-enters only in S+1.

%==============================================================================
\section{Staging}
\label{sec:staging}

\begin{description}[leftmargin=2.2em,itemsep=2pt]
  \item[S+0.] \(M_0\) on Tier 1, then Tier 2, with the clock triage.  Gate: P1 paired CI, P2 memory,
  and Pareto.  Current executable: \code{experiments/p1\_synthetic.py --variant m0}, compared
  against \code{--variant control} and \code{--variant mclk}.
  \item[S+1, P3.] A separate \(M_{\rm MoE}\) ablation uses plain or sparse routing and
  marginal-load balance only.  Dirichlet--UCB stays deleted from the science path.  Gate after S+0
  passes: \(\Delta_{\rm MoE}>0\) at matched budget; \(\mathrm{AEBR}\), route-depth NMI, and expert
  utilization are mechanism diagnostics.  This stage is intentionally not wired into
  \code{experiments/p1\_synthetic.py}.
  \item[S+2, KV axis.] Terminal hidden-state cache (emitted in log contracts as
  \code{hidden\_state\_proxy\_not\_autoregressive\_kv}), with matched training including
  terminal-context attention.  Compare exact multi-depth cache and shared-cache baselines.  Gate:
  \(\alpha_{\rm kv}\approx0\) and quality gap \(\le\tau_{\rm cache}\).  The staged synthetic runner
  emits only a hidden-state memory proxy and explicitly marks the autoregressive KV quality gap as
  untested; an autoregressive LM harness is required before claiming constant-KV correctness.
  \item[S+3, deployment.] Add int6, MLA/GQA, sparse inference, and other deployment features.
  Gate: post-quant Pareto with quantization gap reported.  The current staged runner implements a
  fake-int6 linear-weight sensitivity check only; it is outside the P1 harness.
  \item[Fallback.] Fixed-point pressure only if S+0 fails by instability, severe overthinking, or
  extrapolation.  Use the weakest sufficient stabilizer first.
\end{description}

%==============================================================================
\section{Pipeline Simplification Audit}
\label{sec:pipeline-audit}

The active pipeline should have the smallest interface that can falsify or support P1.  A mechanism
is included only if removing it would make the P1/P2 question unanswerable.  This gives the
following first-principles pruning rule:
\[
  \text{keep}(m)
  \Longleftrightarrow
  m\ \text{is necessary to test P1/P2 now, or is the positive-control/clock diagnostic needed to localize failure.}
\]
Everything else is staged.

\begin{center}
\footnotesize
\renewcommand{\arraystretch}{1.18}
\begin{tabularx}{\textwidth}{@{}L{0.28\textwidth} L{0.34\textwidth} L{0.28\textwidth}@{}}
\toprule
\textbf{Pipeline surface} & \textbf{First-principles critique} & \textbf{Current fix} \\
\midrule
MoE, static-route MoE, expert-basis metrics & P3 mechanisms cannot rescue a failed P1 result and make failure attribution ambiguous. & Removed from the P1 harness; available only in the separate S+1 diagnostic runner. \\
Dirichlet--UCB, MLA, MoS, cache modes & Deployment/capacity features do not answer whether recurrence itself buys utility. & Dirichlet--UCB/MLA/MoS remain excluded; cache appears only as an S+2 proxy diagnostic. \\
Mean-only \(G_T\) reporting & A tiny positive mean can be noise and should not be promoted. & Pair rows now include \(G_{\rm nll}\), \(G_{\rm acc}\), and confidence intervals. \\
Hardcoded task budget in the shell runner & A failed positive control could be task hardness, insufficient iterations, or model failure; hardcoding prevents localization. & Runner exposes only task and budget knobs such as \code{SEQ\_LEN}, \code{ITERATIONS}, \code{TRAIN\_DEPTHS}, and \code{PAIRS}. \\
Legacy hypothesis queue & It still described the old RevDEQ/Parcae/MoE stack as active, steering agents toward refuted or non-P1 work. & The active snapshot is final-minimal P1; the older dense-MoE ledger is historical evidence. \\
Soft Lyapunov or Lipschitz-band pressure & It optimizes a diagnostic proxy rather than task utility and previously hurt validation. & Launch validation rejects nonzero \code{lyapunov\_coef}; FP quantities remain diagnostics. \\
\bottomrule
\end{tabularx}
\end{center}

This audit leaves one intentionally small P1 promotion surface:
\[
  \code{control},\quad \code{m0},\quad \code{mclk}.
\]
It has one job: determine whether depth utility is detectable at all.  It should not grow options
for P3, KV, or deployment while the positive-control gate is failing.  The audit also leaves one known
implementation gap: \code{experiments/p1\_synthetic.py} verifies
structural reversibility and P1 detectability, but it does not yet implement custom reversible
activation reconstruction for training memory.  That is intentional staging, not a P2 claim.  P2 is
promoted only after activation-memory slope and the reversible round-trip error
\(\code{\detokenize{recon_rel}}\) (the BPTT gradient-correctness gate) are measured in the actual
memory-saving backward path.

\subsection{Feedback-stage implementation split}

The repository now has two executable surfaces:
\begin{itemize}[leftmargin=1.4em,itemsep=2pt]
  \item \code{experiments/p1\_synthetic.py}: the minimal S+0/P1 harness.  It accepts only
  \code{control}, \code{m0}, and \code{mclk}.
  \item \code{experiments/p1\_feedback\_stages.py}: non-core staged diagnostics.  S+1 runs
  \code{m0}, \code{moe}, and \code{static\_moe}; S+2 emits a hidden-state cache proxy labelled
  \code{hidden\_state\_proxy\_not\_autoregressive\_kv}; S+3 emits fake-int6 K-sweep and pair
  diagnostics.
\end{itemize}
The full feedback-stage command is:
\[
\code{CUDA\_VISIBLE\_DEVICES=6,7\ ITERATIONS=1000\ RUN\_TAG=gpu67\_feedback\_all\_i1000\ bash\ experiments/run\_feedback\_stage\_pipeline.sh}.
\]
This command is useful for executable coverage of the feedback document.  Its S+1/S+2/S+3 rows are
mechanism/deployment diagnostics, not substitutes for a passing S+0 positive-control gate.

%==============================================================================
\section{Current Simplified S+0 DDP Verification}
\label{sec:s0-smoke}

On June 4, 2026, after pruning the executable harness to the necessary P1 surface, the simplified
S+0 pipeline was run with
\[
\code{CUDA\_VISIBLE\_DEVICES=6,7\ ITERATIONS=1000\ RUN\_TAG=gpu67\_s0\_pruned\_i1000}.
\]
The run used two A100 GPUs, \(\ell=16\), train depths \(\{16,32,64\}\), eval depths
\(\{4,8,16,32,64\}\), and 4096 paired evaluation examples per depth pair.

\begin{center}
\scriptsize
\setlength{\tabcolsep}{3pt}
\renewcommand{\arraystretch}{1.12}
\begin{tabularx}{\textwidth}{@{}L{0.13\textwidth} r r r r L{0.22\textwidth} L{0.22\textwidth} r@{}}
\toprule
\textbf{Variant} & \textbf{Params} & \textbf{Peak MB} & \textbf{Best \(K\)} & \textbf{Best loss} & \textbf{\(G_{\rm nll}(16,64)\) [95\% CI]} & \textbf{\(G_{\rm acc}(16,64)\) [95\% CI]} & \textbf{NDR} \\
\midrule
\code{control} & 198{,}521 & 1316.0 & 16 & 4.7930 & \(-0.000477\ [-0.000951,\ -0.000002]\) & \(0.000488\ [-0.000188,\ 0.001165]\) & 0.5122 \\
\code{m0} & 363{,}897 & 2539.1 & 16 & 4.7934 & \(0.000071\ [-0.000270,\ 0.000412]\) & \(-0.000977\ [-0.002767,\ 0.000814]\) & 0.4968 \\
\code{mclk} & 429{,}433 & 2540.1 & 16 & 4.7932 & \(-0.000073\ [-0.000351,\ 0.000205]\) & \(0.000000\ [0.000000,\ 0.000000]\) & 0.5146 \\
\bottomrule
\end{tabularx}
\end{center}

\begin{center}
\scriptsize
\setlength{\tabcolsep}{3pt}
\renewcommand{\arraystretch}{1.12}
\begin{tabularx}{\textwidth}{@{}L{0.15\textwidth} L{0.28\textwidth} L{0.28\textwidth} r r@{}}
\toprule
\textbf{Variant} & \textbf{\(G_{\rm nll}(8,32)\) [95\% CI]} & \textbf{\(G_{\rm acc}(8,32)\) [95\% CI]} & \textbf{NDR} & \textbf{\(R_{\rm rec}\)} \\
\midrule
\code{control} & \(-0.000203\ [-0.001380,\ 0.000974]\) & \(-0.000732\ [-0.001998,\ 0.000534]\) & 0.5151 & N/A \\
\code{m0} & \(0.000203\ [-0.000573,\ 0.000979]\) & \(-0.000244\ [-0.002098,\ 0.001609]\) & 0.5039 & \(9.57{\times}10^{-6}\) \\
\code{mclk} & \(0.000658\ [0.000027,\ 0.001289]\) & \(0.000000\ [0.000000,\ 0.000000]\) & 0.4917 & \(1.12{\times}10^{-5}\) \\
\bottomrule
\end{tabularx}
\end{center}

The verification validates the simplified interface, DDP operation, finite-loss training,
CI-backed pair reporting, and structural reversibility.  It does not promote P1.  The positive
control is negative on the main \(16\to64\) NLL gain and \(\mathrm{NDR}_\epsilon\) remains near one
half.  \(M_{\rm clk}\) shows a small positive \(8\to32\) NLL interval, but the clock diagnostic
cannot rescue the failing positive-control gate.  Following the triage table, the next principled
step is detectability repair: verify the target representation and positive-control capacity, use
a shorter or curriculum-shaped composition task, or add a simpler positive control known to solve
the task before testing \(M_0\).  Adding MoE, cache, quantization, or fixed-point pressure would not
address this blocker.

%==============================================================================
\section{Current Feedback-Stage DDP Verification}
\label{sec:feedback-stage-run}

The separate feedback-stage pipeline was run on June 4, 2026 with
\[
\code{CUDA\_VISIBLE\_DEVICES=6,7\ ITERATIONS=1000\ RUN\_TAG=gpu67\_feedback\_all\_i1000}.
\]
It wrote current S+0 outputs plus S+1/S+2/S+3 diagnostic outputs under
\code{experiments/training\_logs/p1\_feedback\_gpu67\_feedback\_all\_i1000}.  This run verifies
the code coverage requested by the feedback document while preserving the minimal P1 interface.

\begin{center}
\scriptsize
\setlength{\tabcolsep}{3pt}
\renewcommand{\arraystretch}{1.12}
\begin{tabularx}{\textwidth}{@{}L{0.15\textwidth} r r r r L{0.23\textwidth} r r@{}}
\toprule
\textbf{Run} & \textbf{Params} & \textbf{Peak MB} & \textbf{Best \(K\)} & \textbf{Best loss} & \textbf{\(G_{\rm nll}(16,64)\) [95\% CI]} & \textbf{NDR} & \textbf{\(R_{\rm rec}\)} \\
\midrule
\code{s0\_control} & 198{,}521 & 1316.0 & 16 & 4.7930 & \(-0.000477\ [-0.000951,\ -0.000002]\) & 0.5122 & N/A \\
\code{s0\_m0} & 363{,}897 & 2541.0 & 16 & 4.7934 & \(0.000071\ [-0.000270,\ 0.000412]\) & 0.4968 & \(9.57{\times}10^{-6}\) \\
\code{s0\_mclk} & 429{,}433 & 2540.1 & 16 & 4.7932 & \(-0.000073\ [-0.000351,\ 0.000205]\) & 0.5146 & \(1.12{\times}10^{-5}\) \\
\code{s1\_m0} & 363{,}897 & 2539.1 & 16 & 4.7934 & \(0.000071\ [-0.000270,\ 0.000412]\) & 0.4968 & \(9.57{\times}10^{-6}\) \\
\code{s1\_moe} & 958{,}593 & 6142.5 & 16 & 4.7943 & \(-0.000133\ [-0.000396,\ 0.000130]\) & 0.5159 & \(9.47{\times}10^{-6}\) \\
\code{s1\_static} & 958{,}593 & 6142.5 & 16 & 4.7943 & \(-0.000134\ [-0.000393,\ 0.000125]\) & 0.5110 & \(1.05{\times}10^{-5}\) \\
\code{s2\_m0} & 363{,}897 & 2539.1 & 16 & 4.7934 & \(0.000071\ [-0.000270,\ 0.000412]\) & 0.4968 & \(9.57{\times}10^{-6}\) \\
\code{s3\_m0} & 363{,}897 & 2539.1 & 16 & 4.7934 & \(0.000071\ [-0.000270,\ 0.000412]\) & 0.4968 & \(9.57{\times}10^{-6}\) \\
\bottomrule
\end{tabularx}
\end{center}

\begin{center}
\scriptsize
\setlength{\tabcolsep}{4pt}
\renewcommand{\arraystretch}{1.12}
\begin{tabularx}{\textwidth}{@{}L{0.20\textwidth} L{0.70\textwidth}@{}}
\toprule
\textbf{Stage} & \textbf{Diagnostic result} \\
\midrule
S+1 MoE basis & \code{s1\_moe} reported route-depth NMI \(8.02{\times}10^{-6}\), \(\mathrm{AEBR}=1.385\), and expert utilization \(0.9992\).  \code{s1\_static\_moe} reported route-depth NMI \(4.93{\times}10^{-6}\), \(\mathrm{AEBR}=1.512\), and utilization \(0.9995\).  Both are executable mechanism diagnostics, but neither improves the failed S+0 gate. \\
S+2 cache axis & The hidden-state proxy emitted \(\alpha_{\rm exact}=1.0\), \(\alpha_{\rm terminal}=0.0\), and \(\alpha_{\rm shared}=0.0\), with \(2{,}097{,}152\) bytes per terminal hidden state.  The JSON label is \code{hidden\_state\_proxy\_not\_autoregressive\_kv}; the quality gap is explicitly \code{not\_tested\_requires\_autoregressive\_terminal\_context\_attention}. \\
S+3 deployment & Fake-int6 linear-weight sensitivity reported best-loss gap \(+0.000907\).  The post-int6 \(G_{\rm nll}(16,64)\) was \(0.000044\) with 95\% CI \([-0.000295,0.000383]\) and \(\mathrm{NDR}=0.5005\), so quantization does not rescue depth scaling. \\
\bottomrule
\end{tabularx}
\end{center}

The outcome is the desired implementation split.  All feedback stages execute under DDP on the
requested GPU pair, but the S+0 positive-control gate still fails.  Therefore the next scientific
move remains detectability repair, not mechanism escalation.

%==============================================================================
\section{Superseded All-Stage DDP Run}
\label{sec:allstage-run}

Before the final simplification, the full staged pipeline was run on June 4, 2026 with
\[
\code{CUDA\_VISIBLE\_DEVICES=6,7\ STAGES=all\ SEQ\_LEN=8\ ITERATIONS=1000\ RUN\_TAG=gpu67\_allstages\_seq8\_i1000\_fix}.
\]
That now-superseded runner completed and wrote eight stage outputs:
\[
\{\code{s0\_control},\code{s0\_m0},\code{s0\_mclk},
\code{s1\_m0},\code{s1\_moe},\code{s1\_static\_moe},
\code{s2\_m0},\code{s3\_m0}\}.
\]

\begin{center}
\scriptsize
\setlength{\tabcolsep}{3pt}
\renewcommand{\arraystretch}{1.12}
\begin{tabularx}{\textwidth}{@{}L{0.17\textwidth} r r r L{0.22\textwidth} L{0.22\textwidth} r@{}}
\toprule
\textbf{Run} & \textbf{Params} & \textbf{Peak MB} & \textbf{Best loss} & \textbf{\(G_{\rm nll}(16,64)\) [95\% CI]} & \textbf{\(G_{\rm acc}(16,64)\) [95\% CI]} & \textbf{NDR} \\
\midrule
\code{s0\_control} & 197{,}497 & 651.9 & 4.7896 & \(0.000173\ [-0.000449,\ 0.000795]\) & \(0.001465\ [-0.001781,\ 0.004710]\) & 0.4968 \\
\code{s0\_m0} & 362{,}873 & 1244.5 & 4.7899 & \(0.000000\ [-0.000392,\ 0.000392]\) & \(0.000977\ [-0.000938,\ 0.002891]\) & 0.5002 \\
\code{s0\_mclk} & 428{,}409 & 1246.1 & 4.7909 & \(0.000027\ [-0.000331,\ 0.000384]\) & \(0.000000\ [0.000000,\ 0.000000]\) & 0.4949 \\
\code{s1\_moe} & 957{,}569 & 3095.6 & 4.7908 & \(0.000112\ [-0.000228,\ 0.000453]\) & \(-0.000732\ [-0.003027,\ 0.001563]\) & 0.5044 \\
\code{s1\_static} & 956{,}545 & 3095.6 & 4.7914 & \(0.000317\ [-0.000075,\ 0.000709]\) & \(0.000488\ [-0.000865,\ 0.001842]\) & 0.4885 \\
\bottomrule
\end{tabularx}
\end{center}

The longer, easier S+0 run still does not promote P1.  The positive-control lower confidence bound
is negative for both tested depth pairs and \(\mathrm{NDR}_\epsilon\) remains near \(0.5\).  This is
stronger evidence than the 300-iteration smoke: the bottleneck is not a missing MoE, quantizer, or
cache feature.  The next experimental move should be a detectability repair: either make the
synthetic task expose an actually depth-dependent target to the current architecture, increase the
training signal substantially, or add a simpler positive-control architecture that is known to solve
the task before testing \(M_0\).

This run is treated as evidence for deletion, not as an active run shape.  S+1 code-path evidence
was limited to executability: both \code{moe} and \code{static\_moe} trained under DDP and reported
expert utilization \(1.0\), but they did not improve the S+0 gate and are no longer options on the
P1 runner.  S+2 emitted the expected hidden-state proxy slopes:
\[
  \alpha_{\rm exact}=1.0,\qquad
  \alpha_{\rm terminal}=0.0,\qquad
  \alpha_{\rm shared}=0.0,
\]
with \(2{,}097{,}152\) bytes per terminal state at the tested shape.  This remains a synthetic
proxy, not an autoregressive KV-cache proof.  S+3 int6 evaluation ran; the post-int6
\(G_{\rm nll}(16,64)\) for \code{s3\_m0} was \(-0.000046\) with 95\% CI
\([-0.000432,0.000340]\), so deployment quantization also does not rescue depth scaling.

%==============================================================================
\appendix
\section{Prior Diagnostic Evidence: 16x1 Dense-MoE Path}
\label{app:prior}

The previous in-repo path was a \(16\times1\) dense-MoE recurrent model with Dirichlet--UCB routing,
MLA-style expert attention, MoS output head, Parcae-style damping, int6 artifact scoring, and
diagnostic K-sweeps.  It is useful evidence because it failed P1 despite stable fixed-point
diagnostics, motivating the final minimal design.

\begin{center}
\footnotesize
\renewcommand{\arraystretch}{1.12}
\begin{tabularx}{\textwidth}{@{}L{0.35\textwidth} L{0.25\textwidth} L{0.30\textwidth}@{}}
\toprule
\textbf{Metric} & \textbf{Observed value} & \textbf{Interpretation} \\
\midrule
Full validation \bpb{} & 3.702038 & Diagnostic smoke score, not a promoted result. \\
\(G_T(16,64)\), \(G_T(16,128)\), \(G_T(32,128)\) & \(-0.012196\), \(-0.012865\), \(-0.004076\) & Extra recurrence hurt paired task score. \\
\(\mathrm{NDR}_\epsilon\) & 1.000000 & Every tested pair degraded under the chosen threshold. \\
\(\rho_F\), \(\sigma_{\max}(J_F)\) at \(K=128\) & 0.7525, 0.9633 & Stable/cache-ready diagnostics did not imply useful recurrence. \\
Route-depth NMI and utilization & near 0; near 0.99 & Experts were broadly used but the mixture was nearly static across depth. \\
Peak VRAM & 35632 MB & Memory feasible for the smoke profile. \\
\bottomrule
\end{tabularx}
\end{center}

\begin{center}
\scriptsize
\setlength{\tabcolsep}{4pt}
\renewcommand{\arraystretch}{1.10}
\begin{tabular}{@{}r r r r r r r r r r@{}}
\toprule
\(K\) & \bpb{} & ED$_u$ & ED$_\ell$ & NMI mean & NMI max & util & \(\rho_F\) & \(\sigma_{\max}\) & iter rel. \\
\midrule
4   & 3.7478 & 1.0484   & 1.0400  & 0.0001 & 0.0002 & 0.9944 & 0.7645 & 0.9853 & 0.1193 \\
8   & 3.7556 & 1.1531   & 1.1738  & 0.0001 & 0.0001 & 0.9930 & 0.7564 & 0.9632 & 0.0317 \\
16  & 3.7625 & 1.5920   & 2.3193  & 0.0001 & 0.0001 & 0.9919 & 0.7537 & 0.9623 & 0.0061 \\
24  & 3.7655 & 2.7455   & 5.7373  & 0.0000 & 0.0001 & 0.9915 & 0.7530 & 0.9284 & 0.0027 \\
32  & 3.7670 & 5.3277   & 10.0747 & 0.0000 & 0.0000 & 0.9913 & 0.7529 & 0.9533 & 0.0018 \\
64  & 3.7687 & 24.0905  & 29.1213 & 0.0000 & 0.0000 & 0.9909 & 0.7528 & 0.9169 & 0.0014 \\
128 & 3.7690 & 69.4624  & 62.6883 & 0.0000 & 0.0000 & 0.9908 & 0.7525 & 0.9633 & 0.0013 \\
192 & 3.7690 & 113.3975 & 88.6915 & 0.0000 & 0.0000 & 0.9907 & 0.7526 & 1.0148 & 0.0013 \\
\bottomrule
\end{tabular}
\end{center}

The root-cause read is objective and measurement mismatch: the old stack contained many mechanisms,
but no depth-hard positive control and no clean way to localize whether failure came from the
harness, autonomous depth signal, block expressivity, or MoE routing.  The final minimal protocol
fixes that attribution problem before revisiting MoE or deployment features.

%==============================================================================
\section{Related Work}
\label{sec:related}

The mechanisms retained under the resource/expressiveness goals are each grounded in prior
work; this section records the findings that justify them.

\paragraph{Reversible recurrence and reversible MoE.}
Reversible residual coupling~\cite{revnet} reconstructs layer inputs from outputs, removing
activation storage; Reformer~\cite{reformer} applies it to Transformers. RevFFN~\cite{revffn}
combines reversible two-stream blocks with a Mixture-of-Experts FFN for memory-efficient
full-parameter MoE training (\(\approx\)49\% peak-VRAM reduction, \emph{no} floor/contraction
requirement), directly validating the reversible\(\times\)MoE combination used here. Our
additive coupling (\S\ref{sec:models}) drives each update from the \emph{opposite} stream,
yielding an exact \emph{explicit} inverse (no fixed-point iteration, unlike RevFFN's
cross-attention coupling).

\paragraph{Recurrent depth, weight tying, and expressivity recovery.}
Universal Transformers~\cite{ut}, Huginn~\cite{huginn}, and Parcae~\cite{parcae} treat
recurrent depth as a parameter-saving axis. Weight tying loses per-layer diversity, and sparse
MoE recovers it: MoEUT~\cite{moeut} shows MoE capacity compensates for UT-style layer sharing,
Sparse-Looped-MoE~\cite{sparseloopedmoe} recovers functional diversity across loop passes, and
Relaxed Recursive Transformers~\cite{relaxedrec} relax tied layers with layer-wise LoRA. This
is exactly our mechanism: low-rank experts shared across recurrence steps, composed per step
into high \emph{effective} depth.

\paragraph{Smooth, reversibility-safe sparse routing.}
Hard top-\(k\) routing is non-smooth and reconstruction-unstable under reversibility. ReMoE's
ReLU routing~\cite{remoe} is continuous and fully differentiable yet yields exact-zero (sparse)
gates, with adaptive-\(L_1\) sparsity control and load balancing; it matches top-\(k\) FLOPs,
exceeds top-\(k\) accuracy, and shows 2--3\(\times\) lower activation flip-rate (hence more
reconstruction-stable). We expose \texttt{router\_type}\(\in\{\texttt{relu},\texttt{softmax}\}\)
as a control experiment and, crucially, port ReMoE's \emph{adaptive} mechanism, not merely the
penalty: a fixed \(L_1\) penalty collapses the router (every weight \(\to0\)), whereas the
adaptive \(\lambda_{\rm route}\) controller (\S\ref{sec:models}) targets a sparsity level and
holds \(\code{\detokenize{active_frac}}\) near it, paired with the per-expert load-balanced aux.

\paragraph{Effective-depth and rank metrics.}
The recurrence-equivalence exponent \(\varphi\)~\cite{isodepth} is \emph{the} recurrence-effectiveness
metric: it measures whether looping a block \(r\) times yields the capacity of \(r\) unique blocks
(\(\varphi{=}1\)) or none (\(\varphi{=}0\)), separating true loop capacity from token-budget effects;
truncated BPTT lowers \(\varphi\), motivating our full reversible BPTT (paper reference
\(\varphi\!\approx\!0.46\)).  We estimate \(\varphi\) as a \emph{fitted parameter} of the Iso-Depth
joint scaling law from a pretraining sweep over recurrence counts \(r\) (looped models only --- the
split below already encodes ``loops vs unique blocks'', so no untied baseline is needed).  At fixed
data \(D\) the \(B D^{-\beta}\) term is constant and absorbed into the irreducible offset \(E\):
\begin{equation}
L(r) \;=\; E \;+\; A\,\bigl(N_{\mathrm{once}} + r^{\varphi}\,N_{\mathrm{rec}}\bigr)^{-\alpha},
\label{eq:isodepth}
\end{equation}
where \(N_{\mathrm{once}}\) is the count of non-recurrent parameters (token/positional embeddings,
readout norm and the MoS head --- the tied output embedding counted once) applied once per forward,
and \(N_{\mathrm{rec}}\) is the count of recurrent (looped-block) parameters (the reversible \(F/G\)
delta blocks plus the step embedding).  The split \(N_{\mathrm{once}} + r^{\varphi} N_{\mathrm{rec}}\)
is what encodes the question: \(\varphi{=}1\) means looping \(r\) times buys the capacity of \(r\)
unique blocks, \(\varphi{=}0\) means looping buys nothing (loss flat in \(r\)).  We fit
\((\varphi,\alpha,E,A)\) with \texttt{scipy.optimize.curve\_fit} (bounds \(\varphi\!\in\![0,1]\),
\(\alpha\!\in\!(0,3]\), \(A>0\), \(E\!\in\![0,\min_r L]\); multi-start over \((\varphi_0,\alpha_0)\) for
robustness) to the per-\(r\) trained validation loss collected by \code{experiments/measure\_phi.py}
(an \emph{expensive} multi-train experiment that retrains the model once per \(r\) at matched
data/steps/seed, fixing recurrence depth via a single-element \code{--k-set} and disabling the
no-degradation hinge so each model is a clean pure-depth-\(r\) train).  This is distinct from the
cheap per-run \emph{eval-depth} proxy \code{phi\_eval} (\S\ref{sec:protocol}), which is only an OLS
log-slope of \emph{one} trained model's loss versus \emph{inference} depth \(K\): \code{phi\_eval}
gauges test-time depth utilization, whereas the train-\(r\) \(\varphi\) of Eq.~\eqref{eq:isodepth}
(\code{phi\_isodepth}) measures train-time loop capacity and is the principled metric.  With only
\(\sim\!5\) values of \(r\) the four-parameter fit is small-sample and weakly identified
(\(\varphi/\alpha/A\) partially trade off), so it recovers \(\varphi\) near-exactly on clean data but
is noise-fragile, especially at low \(\varphi\); reporting tight confidence intervals needs more \(r\)
and/or repeated seeds.  Effective rank (spectral entropy of singular values)~\cite{effrank} sets each
low-rank matrix's ``sufficient rank'' on the resource/expressiveness frontier.

\paragraph{Attention and output expressiveness under a resource budget.}
Multi-head Latent Attention~\cite{deepseekv2} low-rank-compresses the KV cache (to 4--14\% of
MHA), the basis for the MLA block. The Mixture-of-Softmaxes head~\cite{mos} breaks the
softmax-rank bottleneck for output-distribution expressiveness; at vocabulary size 1024 its
per-component cost is small.

\begin{thebibliography}{99}

\bibitem{barrington}
D. A. M. Barrington.
\newblock Bounded-width polynomial-size branching programs recognize exactly those languages in
\(\mathrm{NC}^{1}\).
\newblock \emph{Journal of Computer and System Sciences}, 38(1):150--164, 1989.
\newblock \url{https://doi.org/10.1016/0022-0000(89)90037-8}.

\bibitem{ut}
M. Dehghani, S. Gouws, O. Vinyals, J. Uszkoreit, and L. Kaiser.
\newblock Universal Transformers.
\newblock \emph{International Conference on Learning Representations}, 2019.
\newblock \url{https://arxiv.org/abs/1807.03819}.

\bibitem{huginn}
J. Geiping, S. McLeish, N. Jain, J. Kirchenbauer, S. Singh, B. R. Bartoldson,
B. Kailkhura, A. Bhatele, and T. Goldstein.
\newblock Scaling up Test-Time Compute with Latent Reasoning: A Recurrent Depth Approach.
\newblock arXiv:2502.05171, 2025.
\newblock \url{https://arxiv.org/abs/2502.05171}.

\bibitem{isodepth}
K. Schwethelm, D. Rueckert, and G. Kaissis.
\newblock How Much Is One Recurrence Worth? Iso-Depth Scaling Laws for Looped Language Models.
\newblock arXiv:2604.21106, 2026.
\newblock \url{https://arxiv.org/abs/2604.21106}.

\bibitem{hyperconnections}
D. Zhu, H. Huang, Z. Huang, Y. Zeng, Y. Mao, B. Wu, Q. Min, and X. Zhou.
\newblock Hyper-Connections.
\newblock arXiv:2409.19606, 2024.
\newblock \url{https://arxiv.org/abs/2409.19606}.

\bibitem{mhc}
DeepSeek-AI.
\newblock mHC: Manifold-Constrained Hyper-Connections (doubly-stochastic / Birkhoff-polytope mixing
via Sinkhorn--Knopp for stable, signal-conserving residual mixing).
\newblock arXiv:2512.24880, 2025.
\newblock \url{https://arxiv.org/abs/2512.24880}.

\bibitem{jpmhc}
\newblock Dynamical Isometry via Orthogonal Hyper-Connections (strictly-invertible orthogonal mixing
matrix; the reversibility-native variant).
\newblock arXiv:2602.18308, 2026.
\newblock \url{https://arxiv.org/abs/2602.18308}.

\bibitem{cot}
Z. Li, H. Liu, D. Zhou, and T. Ma.
\newblock Chain of Thought Empowers Transformers to Solve Inherently Serial Problems.
\newblock \emph{International Conference on Learning Representations}, 2024.
\newblock \url{https://arxiv.org/abs/2402.12875}.

\bibitem{melt}
V. C. Vendrell, A. P. Masdemont, N. Grillo, J. Ros-Giralt, A. Behboodi, and
F. V. Massoli.
\newblock Memory-Efficient Looped Transformer: Decoupling Compute from Memory in Looped Language Models.
\newblock arXiv:2605.07721, 2026.
\newblock \url{https://arxiv.org/abs/2605.07721}.

\bibitem{moeut}
R. Csordas, K. Irie, J. Schmidhuber, C. Potts, and C. D. Manning.
\newblock MoEUT: Mixture-of-Experts Universal Transformers.
\newblock \emph{NeurIPS}, 2024.
\newblock \url{https://arxiv.org/abs/2405.16039}.

\bibitem{revnet}
A. N. Gomez, M. Ren, R. Urtasun, and R. B. Grosse.
\newblock The Reversible Residual Network: Backpropagation Without Storing Activations.
\newblock \emph{NeurIPS}, 2017.
\newblock \url{https://arxiv.org/abs/1707.04585}.

\bibitem{reformer}
N. Kitaev, {\L}. Kaiser, and A. Levskaya.
\newblock Reformer: The Efficient Transformer.
\newblock \emph{International Conference on Learning Representations}, 2020.
\newblock \url{https://arxiv.org/abs/2001.04451}.

\bibitem{revffn}
RevFFN: Memory-Efficient Full-Parameter Fine-Tuning of Mixture-of-Experts LLMs with Reversible Blocks.
\newblock arXiv:2512.20920, 2025.
\newblock \url{https://arxiv.org/abs/2512.20920}.

\bibitem{remoe}
Z. Wang, J. Chen, and J. Zhu.
\newblock ReMoE: Fully Differentiable Mixture-of-Experts with ReLU Routing.
\newblock \emph{International Conference on Learning Representations}, 2025.
\newblock \url{https://arxiv.org/abs/2412.14711}.

\bibitem{switch}
W. Fedus, B. Zoph, and N. Shazeer.
\newblock Switch Transformers: Scaling to Trillion Parameter Models with Simple and Efficient Sparsity.
\newblock \emph{Journal of Machine Learning Research}, 23(120):1--39, 2022.
\newblock \url{https://arxiv.org/abs/2101.03961}.

\bibitem{deepseekv3}
DeepSeek-AI.
\newblock DeepSeek-V3 Technical Report (auxiliary-loss-free load balancing).
\newblock arXiv:2412.19437, 2024.
\newblock \url{https://arxiv.org/abs/2412.19437}.

\bibitem{stmoe}
B. Zoph, I. Bello, S. Kumar, N. Du, Y. Huang, J. Dean, N. Shazeer, and W. Fedus.
\newblock ST-MoE: Designing Stable and Transferable Sparse Expert Models (router z-loss).
\newblock arXiv:2202.08906, 2022.
\newblock \url{https://arxiv.org/abs/2202.08906}.

\bibitem{relaxedrec}
S. Bae, A. Fisch, H. Harutyunyan, Z. Ji, S. Kim, and T. Schuster.
\newblock Relaxed Recursive Transformers: Effective Parameter Sharing with Layer-wise LoRA.
\newblock arXiv:2410.20672, 2024.
\newblock \url{https://arxiv.org/abs/2410.20672}.

\bibitem{sparseloopedmoe}
Sparse Looped-MoE: Recovering Functional Diversity Across Loop Passes.
\newblock arXiv:2605.09165, 2026.
\newblock \url{https://arxiv.org/abs/2605.09165}.

\bibitem{parcae}
Parcae: Stable Looped Language Models with Loop Depth as a Scaling Axis.
\newblock arXiv:2604.12946, 2026.
\newblock \url{https://arxiv.org/abs/2604.12946}.

\bibitem{deepseekv2}
DeepSeek-AI.
\newblock DeepSeek-V2: A Strong, Economical, and Efficient Mixture-of-Experts Language Model.
\newblock arXiv:2405.04434, 2024.
\newblock \url{https://arxiv.org/abs/2405.04434}.

\bibitem{mos}
Z. Yang, Z. Dai, R. Salakhutdinov, and W. W. Cohen.
\newblock Breaking the Softmax Bottleneck: A High-Rank RNN Language Model.
\newblock \emph{International Conference on Learning Representations}, 2018.
\newblock \url{https://arxiv.org/abs/1711.03953}.

\bibitem{effrank}
O. Roy and M. Vetterli.
\newblock The Effective Rank: A Measure of Effective Dimensionality.
\newblock \emph{European Signal Processing Conference (EUSIPCO)}, 2007.

\bibitem{vinfo}
Y. Xu, S. Zhao, J. Song, R. Stewart, and S. Ermon.
\newblock A Theory of Usable Information Under Computational Constraints.
\newblock \emph{International Conference on Learning Representations}, 2020.
\newblock \url{https://arxiv.org/abs/2002.10689}.

\bibitem{vusable}
K. Ethayarajh, Y. Choi, and S. Swayamdipta.
\newblock Understanding Dataset Difficulty with \(\mathcal{V}\)-Usable Information.
\newblock \emph{International Conference on Machine Learning}, 2022.
\newblock \url{https://arxiv.org/abs/2110.08420}.

\end{thebibliography}

\end{document}
